% ============================================================
% The Elements of Computational Worlds
% Building a 4X Civilization Simulator from First Principles
% in Spherepop
%
% Flyxion
% ============================================================

\documentclass[11pt,a4paper,openany]{book}

% ── Packages ────────────────────────────────────────────────
\usepackage{iftex}
\ifluatex
  % lualatex: use fontspec for proper small caps
  \usepackage{fontspec}
  \setmainfont{Latin Modern Roman}[SmallCapsFont={Latin Modern Roman Caps}]
  \setmonofont{Latin Modern Mono}[Scale=0.9]
\else
  % pdflatex: standard OT1 fonts
\fi
\usepackage{amsmath,amssymb,amsthm}
\usepackage{mathtools}
\usepackage{geometry}
\usepackage{hyperref}
\usepackage[dvipsnames,svgnames,x11names]{xcolor}
\usepackage{listings}
\usepackage{tcolorbox}
\tcbuselibrary{breakable,skins}
\usepackage{booktabs}
\usepackage{array}
\usepackage{enumitem}
\usepackage{tikz}
\usepackage{fancyhdr}
\usepackage{titlesec}
\usepackage{epigraph}
\usepackage{mdframed}
\usepackage{multicol}
\usepackage{graphicx}
\usepackage{caption}
\usepackage{subcaption}
\usepackage{xspace}

\usetikzlibrary{arrows.meta, positioning, shapes.geometric, fit, backgrounds,
                calc, matrix, decorations.pathreplacing}

% ── Page geometry ───────────────────────────────────────────
\geometry{
  a4paper,
  top=1.1in, bottom=1.1in,
  inner=1.25in, outer=1.0in,
  marginparwidth=0.7in,
  headsep=0.25in
}

% ── Colors ──────────────────────────────────────────────────
\definecolor{sphereblue}{RGB}{30, 80, 160}
\definecolor{sphereteal}{RGB}{20, 140, 140}
\definecolor{sphereamber}{RGB}{210, 130, 20}
\definecolor{spherered}{RGB}{180, 40, 40}
\definecolor{spheregray}{RGB}{90, 90, 100}
\definecolor{codebg}{RGB}{245, 247, 250}
\definecolor{codeborder}{RGB}{180, 190, 210}
\definecolor{projectbg}{RGB}{240, 248, 255}
\definecolor{defbg}{RGB}{250, 248, 240}

% ── Hyperref ────────────────────────────────────────────────
\hypersetup{
  colorlinks=true,
  linkcolor=sphereblue,
  urlcolor=sphereteal,
  citecolor=spheregray,
  pdftitle={The Elements of Computational Worlds},
  pdfauthor={Flyxion}
}

% ── Listings (Spherepop pseudocode) ─────────────────────────
\lstdefinelanguage{Spherepop}{
  keywords={bubble, gate, inputs, output, logic, if, then, else, and, or, not,
            produce, consume, emit, absorb, fire, latch, counter, tile,
            agent, world, when, while, do, end, true, false, self, let},
  keywordstyle=\color{sphereblue}\bfseries,
  comment=[l]{--},
  commentstyle=\color{spheregray}\itshape,
  stringstyle=\color{sphereteal},
  morestring=[b]",
  sensitive=false,
  literate={->}{{$\to$}}{2} {=>}{{$\Rightarrow$}}{2}
           {AND}{{$\wedge$}}{3} {OR}{{$\vee$}}{3}
           {NOT}{{$\neg$}}{3}
}

\lstset{
  language=Spherepop,
  basicstyle=\ttfamily\small,
  backgroundcolor=\color{codebg},
  frame=single,
  rulecolor=\color{codeborder},
  framesep=4pt,
  xleftmargin=8pt,
  xrightmargin=8pt,
  breaklines=true,
  tabsize=4,
  showstringspaces=false,
  captionpos=b,
  numbers=left,
  numberstyle=\tiny\color{spheregray},
  numbersep=6pt
}

% ── Theorem environments ─────────────────────────────────────
\theoremstyle{plain}
\newtheorem{theorem}{Theorem}[chapter]
\newtheorem{lemma}[theorem]{Lemma}
\newtheorem{proposition}[theorem]{Proposition}
\newtheorem{corollary}[theorem]{Corollary}

\theoremstyle{definition}
\newtheorem{definition}[theorem]{Definition}
\newtheorem{example}[theorem]{Example}
\newtheorem{construction}[theorem]{Construction}
\newtheorem{axiom}[theorem]{Axiom}
\newtheorem*{gamerul}{Rule}

\theoremstyle{remark}
\newtheorem{remark}[theorem]{Remark}
\newtheorem{note}[theorem]{Note}
\newtheorem{observation}[theorem]{Observation}

% ── Custom environments ──────────────────────────────────────

% Background box — mdframed to avoid tcolorbox font-state bleed
\newenvironment{background}{%
  \begin{mdframed}[
    linecolor=sphereblue!40, backgroundcolor=sphereblue!5,
    linewidth=1pt, innertopmargin=8pt, innerbottommargin=8pt,
    innerleftmargin=10pt, innerrightmargin=10pt,
    skipabove=10pt, skipbelow=6pt
  ]%
  \setlength{\parindent}{0pt}%
  \noindent{\textcolor{sphereblue}{\textbf{Background}}}\par\smallskip\ignorespaces
}{\end{mdframed}}

% Specification box
\newenvironment{specification}{%
  \begin{mdframed}[
    linecolor=sphereteal!60, backgroundcolor=sphereteal!5,
    linewidth=1pt, innertopmargin=8pt, innerbottommargin=8pt,
    innerleftmargin=10pt, innerrightmargin=10pt,
    skipabove=10pt, skipbelow=6pt
  ]%
  \setlength{\parindent}{0pt}%
  \noindent{\textcolor{sphereteal}{\textbf{Specification}}}\par\smallskip\ignorespaces
}{\end{mdframed}}

% Implementation box
\newenvironment{implementation}{%
  \begin{mdframed}[
    linecolor=sphereamber!80, backgroundcolor=sphereamber!5,
    linewidth=1pt, innertopmargin=8pt, innerbottommargin=8pt,
    innerleftmargin=10pt, innerrightmargin=10pt,
    skipabove=10pt, skipbelow=6pt
  ]%
  \setlength{\parindent}{0pt}%
  \noindent{\textcolor{sphereamber!80!black}{\textbf{Implementation}}}\par\smallskip\ignorespaces
}{\end{mdframed}}

% Perspective box
\newenvironment{perspective}{%
  \begin{mdframed}[
    linecolor=spherered!40, backgroundcolor=spherered!4,
    linewidth=1pt, innertopmargin=8pt, innerbottommargin=8pt,
    innerleftmargin=10pt, innerrightmargin=10pt,
    skipabove=10pt, skipbelow=6pt
  ]%
  \setlength{\parindent}{0pt}%
  \noindent{\textcolor{spherered}{\textbf{Perspective}}}\par\smallskip\ignorespaces
}{\end{mdframed}}

% Project box
\newenvironment{project}[1]{%
  \begin{tcolorbox}[
    title={\textbf{Project: #1}},
    colback=projectbg,
    colframe=sphereblue,
    fonttitle=\bfseries\color{white},
    coltitle=white,
    attach boxed title to top left={yshift=-2mm, xshift=4mm},
    boxed title style={colback=sphereblue, sharp corners},
    sharp corners,
    breakable,
    before upper={\normalfont\setlength{\parindent}{0pt}\tolerance=9999\hbadness=9999}
  ]
}{\end{tcolorbox}}

% Definition box
\newenvironment{defbox}[1]{%
  \begin{tcolorbox}[
    title={\textbf{Definition:} #1},
    colback=defbg,
    colframe=sphereamber!80,
    fonttitle=\bfseries,
    coltitle=sphereamber!60!black,
    sharp corners,
    breakable,
    before upper={\normalfont\setlength{\parindent}{0pt}\tolerance=9999\hbadness=9999}
  ]
}{\end{tcolorbox}}

% World mechanic callout
\newenvironment{mechanic}[1]{%
  \begin{tcolorbox}[
    title={\textbf{World Mechanic:} #1},
    colback=sphereteal!5,
    colframe=sphereteal,
    fonttitle=\bfseries\color{white},
    coltitle=white,
    sharp corners,
    breakable,
    before upper={\normalfont\setlength{\parindent}{0pt}\tolerance=9999\hbadness=9999}
  ]
}{\end{tcolorbox}}

% ── Headers/footers ──────────────────────────────────────────
\pagestyle{fancy}
\fancyhf{}
\fancyhead[LE]{\small\textcolor{spheregray}{\textit{\leftmark}}}
\fancyhead[RO]{\small\textcolor{spheregray}{\textit{\rightmark}}}
\fancyfoot[C]{\small\textcolor{spheregray}{\thepage}}
\renewcommand{\headrulewidth}{0.3pt}

% ── Chapter title style ──────────────────────────────────────
\titleformat{\chapter}[display]
  {\normalfont\huge\bfseries\color{sphereblue}}
  {\normalfont\large\color{spheregray}\chaptername\ \thechapter}
  {8pt}
  {\Huge}
  [\vspace{2pt}\color{sphereblue!30}\rule{\linewidth}{1.5pt}]

\titlespacing{\chapter}{0pt}{20pt}{30pt}

\titleformat{\section}
  {\normalfont\Large\bfseries\color{sphereblue!80!black}}
  {\thesection}{1em}{}

\titleformat{\subsection}
  {\normalfont\large\bfseries\color{sphereteal}}
  {\thesubsection}{1em}{}

% ── Epigraph style ───────────────────────────────────────────
\setlength{\epigraphwidth}{0.65\linewidth}
\renewcommand{\epigraphflush}{flushright}

% ── Misc macros ──────────────────────────────────────────────
\newcommand{\Sp}{\textsc{Spherepop}\xspace}
\ifluatex
  \newcommand{\Spbox}{\textsc{Spherepop}\xspace}
\else
  \newcommand{\Spbox}{\textbf{Spherepop}\xspace}
\fi
\newcommand{\bool}{\mathbb{B}}
\newcommand{\AND}{\ensuremath{\wedge}}
\newcommand{\OR}{\ensuremath{\vee}}
\newcommand{\NOT}{\ensuremath{\neg}}
\newcommand{\XOR}{\ensuremath{\oplus}}
\newcommand{\nand}{\ensuremath{\barwedge}}
\newcommand{\tile}[1]{\mathtt{T}_{#1}}
\newcommand{\agent}[1]{\mathcal{A}_{#1}}
\newcommand{\world}{\ensuremath{\mathcal{W}}}
\newcommand{\state}{\ensuremath{\mathcal{S}}}
\newcommand{\trans}[1]{\xrightarrow{#1}}
\newcommand{\latch}[1]{\mathtt{Q}_{#1}}
\newcommand{\signal}[1]{\sigma_{#1}}

% World-transition operator
\newcommand{\Fworld}{F_{\mathcal{W}}}
\newcommand{\Feco}{F_{\mathrm{eco}}}
\newcommand{\Fagent}{F_{\mathrm{agent}}}
\newcommand{\Fecon}{F_{\mathrm{econ}}}
\newcommand{\Fcombat}{F_{\mathrm{combat}}}
\newcommand{\Fdip}{F_{\mathrm{dip}}}

% Fixed-point shorthand
\newcommand{\fx}[1]{\langle #1 \rangle_{\mathrm{fx}}}

% Chapter quote (typographically stable replacement for epigraph in chapter openers)
\newcommand{\chapterquote}[2]{%
  \par\vspace{0.8em}%
  \begin{flushright}%
    \begin{minipage}{0.65\textwidth}%
      {\itshape #1\par}%
      \ifx&#2&\else\vspace{0.3em}{\raggedleft\small--- #2\par}\fi%
    \end{minipage}%
  \end{flushright}%
  \vspace{0.8em}\par\noindent}

% mdframed global defaults (prevent bad page splits)
\mdfsetup{skipabove=10pt, skipbelow=8pt}

% Hyphenation exceptions
\hyphenation{Sphere-pop}

% ── Document ─────────────────────────────────────────────────
\begin{document}

\sloppy
\setlength{\emergencystretch}{3em}
\hfuzz=4pt
\frontmatter

% ── Title page ───────────────────────────────────────────────
\begin{titlepage}
  \centering
  \vspace*{1.5cm}
  {\color{sphereblue}\rule{\linewidth}{2pt}}\\[0.6cm]
  {\Huge\bfseries\color{sphereblue} The Elements of\\[0.3em]
   Computational Worlds}\\[0.5cm]
  {\color{sphereblue!30}\rule{\linewidth}{1pt}}\\[0.4cm]
  {\Large\itshape\color{sphereteal}
   Building a 4X Civilization Simulator\\
   from First Principles in \textsc{Spherepop}}\\[1.2cm]
  {\color{sphereblue!30}\rule{0.4\linewidth}{0.4pt}}\\[0.5cm]
  {\large Flyxion}\\[0.2cm]
  {\small\color{spheregray} Independent Researcher}\\[2.5cm]
  \begin{tikzpicture}[scale=1.0]
    % Schematic of the abstraction tower
    \foreach \y/\col/\lbl in {
      0/spheregray!60/{\small Signals},
      1/sphereblue!40/{\small Logic Gates},
      2/sphereblue!60/{\small Memory \& State},
      3/sphereteal!60/{\small Tiles \& Terrain},
      4/sphereteal!80/{\small Agents \& Economies},
      5/sphereblue!80/{\small Civilization \& History}
    }{
      \fill[\col, rounded corners=3pt]
        (-3.2+\y*0.12, \y*0.65) rectangle (3.2-\y*0.12, \y*0.65+0.55);
      \node at (0, \y*0.65+0.27) {\lbl};
    }
    \draw[-Stealth, thick, sphereblue] (0,-0.15) -- (0,-0.55)
      node[below]{\small\itshape\color{spheregray} from gates to worlds};
  \end{tikzpicture}
  \vfill
  {\color{sphereblue!30}\rule{\linewidth}{1pt}}\\[0.3cm]
  {\small\color{spheregray} First Edition}
\end{titlepage}

% ── Copyright ────────────────────────────────────────────────
\thispagestyle{empty}
\vspace*{\fill}
\noindent\textit{The Elements of Computational Worlds}\\
\noindent Copyright \copyright\ Cogito Ergo Sum.\\[0.5em]
\noindent All rights reserved. No part of this work may be reproduced
without prior written permission of the author.\\[1em]
\noindent This textbook was developed in conjunction with the
\Sp{} computational framework, the RSVP field-theoretic programme,
and the KES (Kinetic-Event Synthesis) formalism.\\[1em]
\noindent\textit{``A world is a layered state machine whose global
history emerges from local logical permissions.''}
\vspace*{\fill}
\newpage

% ── Preface ──────────────────────────────────────────────────
\chapter*{Preface}
\addcontentsline{toc}{chapter}{Preface}

\chapterquote{The elements of every science are at once its lowest and its
  highest truths: lowest because they are common to all, highest because
  they are the key to all.
}{--- paraphrase of Aristotle, \textit{Prior Analytics}}

This book is about building worlds from scratch. Not artistic worlds, not
narrative worlds, but \emph{computational} worlds: persistent, rule-governed
systems of state that support exploration, expansion, exploitation, and
conflict---the four pillars of the 4X strategy genre.

The organizing thesis is simple. A 4X world is a civilization-scale circuit.
Everything that makes such a world interesting---territory, resources, memory,
agency, diplomacy, history---emerges from a small set of computational
primitives: signals, logic gates, latches, counters, and state machines.

This book follows the pedagogical architecture of Noam Nisan and Shimon
Schocken's masterwork \textit{The Elements of Computing Systems}. Each
chapter introduces an abstraction, provides a formal specification, works
through an implementation in the \Sp{} language, offers a philosophical
perspective, and closes with a construction project. The ascent of the
book mirrors the ascent of complexity:
\[
  \text{Signals} \to \text{Gates} \to \text{Memory}
  \to \text{Tiles} \to \text{Agents}
  \to \text{Economies} \to \text{Diplomacy}
  \to \text{History}
\]
What distinguishes this from a standard computing systems course is the
endpoint. Nisan and Schocken build a general-purpose computer. This book
builds a world. Every abstraction is motivated not by the needs of an
operating system, but by the demands of a civilization that must remember,
perceive, decide, and endure.

\Sp{} is our medium throughout. It treats computation as the interaction of
\emph{bubbles}: self-contained units that may carry values, execute rules,
spawn children, consume neighbors, or pass signals along. Bubbles are
simultaneously data, process, interface, and environment. This unification
makes \Sp{} unusually well suited to world simulation: a tile is a bubble.
An agent is a bubble. A diplomatic treaty is a bubble. The world itself is
a bubble.

The book is divided into five parts:
\begin{enumerate}[leftmargin=*, label=\Roman*.]
  \item \textbf{Foundations of Computation.}
    Signals, Boolean algebra, logic gates, and bubble interaction.
  \item \textbf{State and Memory.}
    Latches, registers, counters, and the persistence of time.
  \item \textbf{Spatial Worlds.}
    Tiles, terrain, visibility, and the geometry of territory.
  \item \textbf{Agents and Economies.}
    Perception, action, resources, trade, and conflict.
  \item \textbf{Civilization and Emergence.}
    Technology, diplomacy, ecology, and the accumulation of history.
\end{enumerate}

A student who completes all five parts and their associated projects will
have constructed, from first principles, a functioning 4X world simulation.
They will also have understood \emph{why} that world works the way it does,
because they will have built every layer themselves.

\bigskip
\noindent\textit{Flyxion}\\
\noindent\textit{Canada}

% ── How to Use This Book ─────────────────────────────────────
\chapter*{How to Use This Book}
\addcontentsline{toc}{chapter}{How to Use This Book}

Each chapter is structured in five sections following the pedagogical
pattern of Nisan and Schocken:

\begin{description}[style=nextline]
  \item[\textcolor{sphereblue}{Background}]
    Motivates the chapter's central concept from first principles.
    Why does a world need this abstraction?

  \item[\textcolor{sphereteal}{Specification}]
    Defines the abstraction precisely. What are its inputs, outputs,
    and behavioral contracts?

  \item[\textcolor{sphereamber!80!black}{Implementation}]
    Constructs the abstraction in \Sp{}. How is it realized
    from the primitives introduced in prior chapters?

  \item[\textcolor{spherered}{Perspective}]
    Reflects on what the abstraction means philosophically, ecologically,
    economically, or strategically.

  \item[\textcolor{sphereblue}{Project}]
    A hands-on construction task. The student builds the abstraction,
    tests it, and integrates it into the growing simulation.
\end{description}

Projects are cumulative. Each project produces a component that later
chapters will use. By the final chapter, the student has assembled a
working civilization simulator from nothing but bubble interactions.

\tableofcontents

\mainmatter

% ============================================================
% PART I: FOUNDATIONS OF COMPUTATION
% ============================================================

\part{Foundations of Computation}

\chapterquote{Before a world can have empires, cities, or history, it must
  have a logic of difference.
}{}

% ── Chapter 1 ────────────────────────────────────────────────
\chapter{Boolean Worlds}
\label{ch:boolean}

\chapterquote{Something is passable or impassable, known or unknown,
  owned or unowned. These distinctions are prior to everything else.
}{}

\begin{background}
  A 4X world is, at its core, a system of distinctions. Before there can
  be empires or wars, there must be a way to say: this tile is mine and
  that one is not; this unit is alive and that one is dead; this technology
  is known and that one is not. Such distinctions are not merely
  descriptive---they drive computation. The world \emph{acts} differently
  depending on whether a tile is owned, whether a unit is visible, whether
  a treaty is active.

  The mathematics of yes/no distinctions is Boolean algebra. Developed by
  George Boole in the 1840s as an algebra of thought, Boolean algebra turns
  out to be the foundation of all digital computation. We begin there.
\end{background}

\section{Signals}

\begin{defbox}{Signal}
  A \emph{signal} is an element of $\bool$. We call $0$ \emph{inactive}
  (or \emph{false}) and $1$ \emph{active} (or \emph{true}).
\end{defbox}

In a 4X context, the same Boolean distinction recurs at every scale:

\begin{center}
\begin{tabular}{lll}
  \toprule
  \textbf{Domain} & \textbf{Signal $= 0$} & \textbf{Signal $= 1$} \\
  \midrule
  Map visibility  & Unexplored           & Explored              \\
  Ownership       & Neutral              & Claimed               \\
  Unit status     & Dead                 & Alive                 \\
  Technology      & Unknown              & Researched            \\
  Terrain         & Impassable           & Passable              \\
  Treaty          & Hostile              & Allied                \\
  Resource node   & Depleted             & Active                \\
  \bottomrule
\end{tabular}
\end{center}

A single bit contains one unit of world-state. Stacked across millions of
tiles and thousands of game-turns, bits become civilizations.

\section{Boolean Operations}

Three operations are primitive:

\begin{definition}[Boolean Operations]
  For $a, b \in \bool$:
  \begin{align}
    a \AND b &= 1 \quad \Leftrightarrow \quad a = 1 \text{ and } b = 1 \\
    a \OR b  &= 1 \quad \Leftrightarrow \quad a = 1 \text{ or } b = 1
                        \text{ (or both)} \\
    \NOT a   &= 1 - a
  \end{align}
\end{definition}

\begin{mechanic}{Permission Logic}
  The most fundamental 4X question is: ``Is this action permitted?''
  Permissions are conjunctions:
  \[
    \texttt{can\_move} = \texttt{passable} \AND \texttt{not\_occupied}
    \AND \texttt{has\_movement\_points}
  \]
  \[
    \texttt{can\_settle} = \texttt{unowned} \AND \texttt{adjacent\_to\_unit}
    \AND \NOT\, \texttt{at\_war}
  \]
  Every game rule is ultimately a Boolean expression over world state.
\end{mechanic}

\section{NAND Universality}

One gate is sufficient for all Boolean computation:

\begin{theorem}[NAND Universality]
  Every Boolean function $f : \bool^n \to \bool^m$ can be expressed
  using only NAND gates, where
  \[
    \text{NAND}(a, b) = \NOT(a \AND b) = 1 - ab.
  \]
\end{theorem}

The derivations:
\begin{align*}
  \NOT(a) &= \text{NAND}(a, a) \\
  a \AND b &= \NOT(\text{NAND}(a,b)) = \text{NAND}(\text{NAND}(a,b),\,\text{NAND}(a,b)) \\
  a \OR b &= \text{NAND}(\text{NAND}(a,a),\,\text{NAND}(b,b))
\end{align*}

This is the first deep fact of the book: \emph{one primitive suffices to build
everything}. In computation, as in worlds, power emerges from iteration of
simple rules.

\begin{specification}
  \textbf{NAND gate.}\par
  \textit{Inputs:}\enspace $a, b \in \bool$.\par
  \textit{Output:}\enspace $\texttt{out} \in \bool$.\par
  \textit{Contract:}\enspace $\texttt{out} = 1$ iff not both $a = 1$ and $b = 1$.
\end{specification}

\begin{implementation}
  In \Spbox , a gate is a named bubble with declared ports and a logic body.
  The NAND gate is the primitive from which all others are derived:

\begin{lstlisting}[caption={NAND gate in \Spbox }]
bubble NAND {
    inputs:  a, b
    output:  out

    -- NAND is the primitive: 0 only when both inputs are 1
    out = not (a and b)
}
\end{lstlisting}

  All subsequent gates in this chapter are specifications, to be implemented
  as compositions of \texttt{NAND} in the Project.
\end{implementation}

\begin{perspective}
  NAND universality is not merely a technical curiosity. It tells us that
  the universe of Boolean functions is reachable from a single operation.
  This is philosophically resonant: a world that begins with only
  ``not-both-true'' can, through composition, express every distinction
  that civilization might require. The architecture of complexity does not
  demand a rich initial vocabulary. It demands iteration.
\end{perspective}

\begin{project}{Boolean Gates}
  \textbf{Objective:} Implement the following gates in \Sp{} using only
  the primitive \texttt{NAND} bubble. Do not use \texttt{and}, \texttt{or},
  or \texttt{not} directly---derive them from \texttt{NAND}.

  \begin{enumerate}
    \item \texttt{NOT(a)} --- logical negation
    \item \texttt{AND(a, b)} --- logical conjunction
    \item \texttt{OR(a, b)} --- logical disjunction
    \item \texttt{XOR(a, b)} --- exclusive disjunction (odd parity)
    \item \texttt{MUX(a, b, sel)} --- multiplexer: output $a$ if
          $\texttt{sel} = 0$, else $b$
    \item \texttt{DMUX(in, sel)} --- demultiplexer: route \texttt{in}
          to output $a$ or $b$ depending on \texttt{sel}
  \end{enumerate}

  For each gate, verify your implementation against its truth table.
  Then implement \texttt{PermissionCheck}, a 4-input bubble that is
  \texttt{true} only when \texttt{passable}, \texttt{visible},
  \texttt{NOT occupied}, and \texttt{NOT at\_war} are all true.
\end{project}

\begin{figure}[htb]
\centering
\begin{tikzpicture}[
  gate/.style={draw, rounded corners=2pt, minimum width=1.3cm,
               minimum height=0.7cm, fill=sphereblue!10, thick},
  wire/.style={thick, color=sphereblue}
]
  \node[gate] (n1) at (0,0) {\small NAND};
  \draw[wire] (-1.8, 0.2) node[left]{\small$a$} -- (n1.west);
  \draw[wire] (-1.8,-0.2) -- (n1.west);
  \draw[wire] (n1.east) -- (2.0,0) node[right]{\small$\neg a$};
  \node[below, color=spheregray, font=\small] at (0,-0.55)
    {NOT $=$ NAND$(a,a)$};
  \begin{scope}[yshift=-2.3cm]
    \node[gate] (n2) at (0,0) {\small NAND};
    \node[gate] (n3) at (2.6,0) {\small NAND};
    \draw[wire] (-1.8, 0.2) node[left]{\small$a$} -- (n2.west);
    \draw[wire] (-1.8,-0.2) node[left]{\small$b$} -- (n2.west);
    \draw[wire] (n2.east) -- (n3.west);
    \draw[wire] (1.55, 0.15) |- (n3.west);
    \draw[wire] (n3.east) -- (4.3,0) node[right]{\small$a \wedge b$};
    \node[below, color=spheregray, font=\small] at (1.3,-0.55)
      {AND $=$ NAND(NAND$(a,b)$, NAND$(a,b)$)};
  \end{scope}
\end{tikzpicture}
\caption{Constructing NOT and AND from NAND: functional completeness in action.}
\label{fig:nand-derivation}
\end{figure}

% ── Chapter 2 ────────────────────────────────────────────────
\chapter{Gates and Bubble Logic}
\label{ch:gates}

\chapterquote{A bubble may contain a value, a rule, a signal, or another
  structure of bubbles.
}{}

\begin{background}
  In the previous chapter we treated gates as abstract functions over
  $\bool$. Now we ask: what is a gate made of, and how does \Spbox  realize
  it? The answer is \emph{bubbles}. A bubble is a self-contained
  computational unit. It has a boundary, a set of ports through which
  signals enter and leave, and an interior that transforms inputs into
  outputs. Crucially, bubbles can contain other bubbles: abstraction is
  nesting.
\end{background}

\section{The Bubble Model}

\begin{defbox}{Bubble}
  A \emph{bubble} $B$ consists of:
  \begin{itemize}
    \item a finite set of \emph{input ports} $I_B$,
    \item a finite set of \emph{output ports} $O_B$,
    \item an \emph{interior} specifying how outputs are computed from inputs,
    \item a \emph{boundary} separating interior from exterior.
  \end{itemize}
  The interior may be primitive (an irreducible gate) or composite
  (an arrangement of sub-bubbles connected by wires).
\end{defbox}

Composing two bubbles $B_1, B_2$ means connecting an output port of $B_1$
to an input port of $B_2$. The resulting composite is itself a bubble.

\section{Multi-Bit Buses}

World state rarely fits in one bit. A tile may have a terrain type coded
in 3 bits, a resource count in 8 bits, a population in 16 bits. We extend
the model to \emph{buses}: ordered tuples of signals.

\begin{definition}[Bus]
  An $n$-bit bus is an element of $\bool^n$. We write buses with subscripts:
  $a[0..n{-}1]$ denotes an $n$-bit value whose $k$-th bit is $a[k]$.
\end{definition}

\begin{defbox}{Half Adder}
  The \emph{half adder} computes the binary sum of two bits:
  \begin{align*}
    \texttt{sum}  &= a \XOR b \\
    \texttt{carry} &= a \AND b
  \end{align*}
\end{defbox}

\begin{defbox}{Full Adder}
  The \emph{full adder} includes a carry-in:
  \begin{align*}
    \texttt{sum}   &= a \XOR b \XOR c \\
    \texttt{carry} &= (a \AND b) \OR (c \AND (a \XOR b))
  \end{align*}
\end{defbox}

Chaining $n$ full adders gives an $n$-bit adder: the circuit that adds
two $n$-bit numbers. Population growth, resource accumulation, and
production totals all reduce to addition.

\section{Bubbles as World Objects}

The insight that will carry us through the entire book is this:
\emph{every world object is a bubble}.

\begin{center}
\begin{tabular}{ll}
  \toprule
  \textbf{World Object} & \textbf{Bubble Type} \\
  \midrule
  Terrain tile          & Primitive state bubble \\
  City                  & Composite bubble (production, population, defense) \\
  Unit                  & Mobile bubble with perception ports \\
  Trade route           & Wire connecting two city bubbles \\
  Treaty                & Shared state bubble between two agent bubbles \\
  Technology tree       & Directed acyclic graph of gate bubbles \\
  History               & Trace of past bubble activations \\
  \bottomrule
\end{tabular}
\end{center}

The world is the outermost bubble. Inside it, civilization emerges from
composition.

\begin{implementation}
\begin{lstlisting}[caption={Multi-bit AND in \Spbox }]
bubble AND16 {
    inputs:  a[0..15], b[0..15]
    output:  out[0..15]

    -- Apply single-bit AND to each bit position
    let i = 0
    while i < 16 do
        out[i] = AND(a[i], b[i])
        i = i + 1
    end
}

bubble Add16 {
    inputs:  a[0..15], b[0..15]
    output:  out[0..15]

    -- Ripple-carry addition across 16 bits
    let carry = 0
    let i = 0
    while i < 16 do
        out[i] = XOR(XOR(a[i], b[i]), carry)
        carry   = OR(AND(a[i], b[i]),
                     AND(carry, XOR(a[i], b[i])))
        i = i + 1
    end
}
\end{lstlisting}
\end{implementation}

\begin{project}{Arithmetic Unit}
  Build the following multi-bit bubbles, each from simpler components
  assembled in prior steps:

  \begin{enumerate}
    \item \texttt{NOT16}: bitwise NOT of a 16-bit bus
    \item \texttt{AND16}: bitwise AND of two 16-bit buses
    \item \texttt{OR16}: bitwise OR of two 16-bit buses
    \item \texttt{MUX16}: 16-bit multiplexer
    \item \texttt{Add16}: 16-bit ripple-carry adder
    \item \texttt{Inc16}: increment a 16-bit value by 1
    \item \texttt{ResourceAdder}: an \texttt{Add16} variant that
          saturates at a maximum value (models bounded resources)
    \item \texttt{PopulationALU}: an arithmetic unit accepting
          \texttt{births}, \texttt{deaths}, and \texttt{migration}
          signals, outputting updated population counts
  \end{enumerate}
\end{project}

\begin{figure}[htb]
\centering
\begin{tikzpicture}[
  gate/.style={draw, minimum width=1.3cm, minimum height=0.75cm,
               fill=sphereteal!10, thick},
  wire/.style={thick, color=sphereteal}
]
  \node[gate] (xor) at (0,1.1)  {\small XOR};
  \node[gate] (and) at (0,-1.1) {\small AND};
  \draw[wire] (-2, 1.35) node[left]{\small$a$} -- (xor.west);
  \draw[wire] (-2, 0.85) -- (xor.west);
  \draw[wire] (-2,-0.85) node[left]{\small$b$} -- (and.west);
  \draw[wire] (-2,-1.35) -- (and.west);
  \draw[wire] (-1.6, 1.35) |- (-0.65, 0.85);
  \draw[wire] (-1.6,-1.35) |- (-0.65,-0.85);
  \draw[wire] (xor.east) -- (2.2, 1.1) node[right]{\small sum};
  \draw[wire] (and.east) -- (2.2,-1.1) node[right]{\small carry};
  \node[below, color=spheregray, font=\small] at (0,-2.0)
    {Half-adder: 1-bit addition with carry};
\end{tikzpicture}
\caption{A half-adder from XOR and AND. This is the first arithmetic circuit: two bits in, sum and carry out.}
\label{fig:half-adder}
\end{figure}


% ============================================================
% PART II: STATE AND MEMORY
% ============================================================

\part{State and Memory}

\chapterquote{Without memory, there is no territory, no ownership,
  no accumulated production, no history, and no strategy.
}{}


\section{Numerical Representation and Fixed-Point Arithmetic}
\label{sec:fixedpoint}

Before worlds can compute resources, growth, or influence, they need numbers
--- not just bits. This section develops the numerical substrate that all
later simulation layers depend on.

\subsection{Unsigned Integers}

An $n$-bit unsigned integer represents a value in $\{0, 1, \ldots, 2^n - 1\}$
via binary positional notation:
\[
  v = \sum_{i=0}^{n-1} b_i \cdot 2^i
\]
where $b_i$ is the $i$-th bit. Addition of two $n$-bit values produces a
$(n+1)$-bit result; the carry bit signals overflow.

\subsection{Fixed-Point Arithmetic for Simulation}

Ecology fields, influence values, and resource yields are continuous in
principle but must be stored in finite registers. We use $k.f$ fixed-point
notation: $k$ bits for the integer part, $f$ bits for the fractional part,
representing values in multiples of $2^{-f}$.

\begin{definition}[Fixed-Point Value]
  A $k.f$ fixed-point register stores the real value
  \[
    v = \frac{\text{raw\_int}}{2^f}
  \]
  where \texttt{raw\_int} is the underlying unsigned integer.
  Addition in fixed-point is ordinary integer addition. Multiplication
  of two $k.f$ values requires a right-shift of $f$ bits to renormalize.
\end{definition}

For the ecological diffusion equation
\[
  x_{t+1} = \alpha x_t + \beta y_t
\]
we store $\alpha, \beta$ as 0.8 fixed-point values (8 fractional bits,
approximating real numbers in $[0,1]$ with resolution $2^{-8} \approx 0.004$).

\begin{defbox}{Saturation Arithmetic}
  To prevent register overflow in resource accumulations, we use
  \emph{saturation arithmetic}: instead of wrapping around, values clamp
  at the maximum (or minimum) representable value.
  \[
    \texttt{sat\_add}(a, b) = \min(a + b,\; 2^n - 1)
  \]
  This is essential for population and resource registers that should never
  go negative or overflow silently.
\end{defbox}

\begin{implementation}
\begin{lstlisting}[caption={Fixed-point multiply in \Spbox }]
bubble FixedMul8 {
    -- Multiply two 8-bit 0.8 fixed-point values
    inputs:  a[8], b[8]
    output:  result[8]

    -- Full 16-bit product, then right-shift 8 to renormalize
    let product[16] = a * b      -- unsigned 16-bit multiply
    result = product[15:8]       -- take upper 8 bits
}
\end{lstlisting}
\end{implementation}

Fixed-point arithmetic is used throughout the book for ecology fields
(Chapter~\ref{ch:ecology}), influence propagation (Chapter~\ref{ch:diplomacy}),
and trajectory entropy calculations (Chapter~\ref{ch:optionality}).


% ── Chapter 3 ────────────────────────────────────────────────
\chapter{Sequential Memory}
\label{ch:memory}

\chapterquote{The past is not the opposite of the present.
  It is the infrastructure of the present.
}{}

\begin{background}
  Everything so far has been \emph{combinational}: given inputs, compute
  outputs immediately, with no memory of prior states. A combinational
  circuit is stateless---it knows only the present moment. But a world
  requires history. A tile must remember whether it has been explored.
  A city must remember its accumulated food surplus. An empire must
  remember its treaties.

  Memory requires \emph{feedback}: a circuit whose output influences its
  future input. The simplest such circuit is the SR latch.
\end{background}

\section{The SR Latch}

\begin{defbox}{SR Latch}
  The \emph{SR latch} has two inputs, \texttt{S} (set) and
  \texttt{R} (reset), and two outputs, \texttt{Q} and $\overline{Q}$.
  Its behavior:
  \begin{align*}
    S = 1, R = 0 &\implies Q := 1 \quad \text{(set)} \\
    S = 0, R = 1 &\implies Q := 0 \quad \text{(reset)} \\
    S = 0, R = 0 &\implies Q \text{ holds its prior value} \quad \text{(memory)} \\
    S = 1, R = 1 &\implies \text{undefined (forbidden)}
  \end{align*}
\end{defbox}

The ``memory'' case is the decisive one. When both inputs are inactive,
the latch \emph{holds}. This is the birth of time in our circuit model.

\begin{mechanic}{Tile Ownership Memory}
  A tile's ownership is implemented as an SR latch:
  \begin{align*}
    S &= \texttt{conquered}(t) \\
    R &= \texttt{liberated}(t) \\
    Q &= \texttt{is\_owned}(t)
  \end{align*}
  When neither conquering nor liberation occurs, the tile remembers its
  owner across turns---which is precisely sovereignty.
\end{mechanic}

\section{The Data Flip-Flop}

The D flip-flop (DFF) is the canonical memory element: it samples its
input \texttt{in} on the rising edge of a clock signal, and holds that
value until the next edge.

\begin{defbox}{Data Flip-Flop (DFF)}
  \begin{align*}
    Q(t+1) &= D(t)
  \end{align*}
  The DFF is the one-tick delay: it stores what its input was one clock
  cycle ago.
\end{defbox}

In \Sp{}, time is a sequence of discrete \emph{ticks}. Each tick
corresponds to one game turn at the world's lowest level.

\section{Registers and RAM}

From DFFs we build registers (multi-bit storage) and RAM arrays
(addressed multi-register storage).

\begin{defbox}{$n$-Bit Register}
  An $n$-bit register stores one $n$-bit word. On each tick, if
  \texttt{load} is active, the register updates to hold \texttt{in};
  otherwise it holds its prior value.
  \[
    Q[n](t+1) = \begin{cases}
      \texttt{in}(t) & \text{if } \texttt{load}(t) = 1 \\
      Q[n](t)        & \text{otherwise}
    \end{cases}
  \]
\end{defbox}

\begin{defbox}{RAM}
  An $n$-word RAM with $k$-bit address width provides:
  \begin{itemize}
    \item \texttt{address[$k$]}: select which of $n = 2^k$ registers to read/write
    \item \texttt{in[$w$]}: data to write
    \item \texttt{load}: write enable
    \item \texttt{out[$w$]}: data read from addressed register
  \end{itemize}
\end{defbox}

\begin{implementation}
\begin{lstlisting}[caption={World-state register in \Spbox }]
bubble TileRegister {
    inputs:  in[16], load
    output:  out[16]

    -- Internal state: one DFF per bit
    latch state[0..15]

    when load = 1 do
        state[0..15] := in[0..15]
    end

    out[0..15] = state[0..15]
}

bubble WorldRAM {
    -- 256 tiles, each holding 16 bits of state
    inputs:  address[8], in[16], load
    output:  out[16]

    latch memory[256][16]

    when load = 1 do
        memory[address] := in
    end

    out = memory[address]
}
\end{lstlisting}
\end{implementation}

\begin{project}{World Memory}
  \begin{enumerate}
    \item Build a \texttt{Bit} bubble (single-bit register) from a DFF.
    \item Build a \texttt{Register16} from 16 \texttt{Bit} bubbles.
    \item Build \texttt{RAM8}, \texttt{RAM64}, \texttt{RAM512},
          \texttt{RAM4K}, \texttt{RAM16K} by recursive composition.
    \item Design a \texttt{TileState} bubble that stores:
          terrain type (3 bits), resource count (8 bits),
          owner ID (4 bits), visibility (1 bit).
    \item Build a \texttt{WorldMap} of $16 \times 16 = 256$ \texttt{TileState}
          bubbles, addressable by $(x, y)$ coordinates.
  \end{enumerate}
\end{project}

\begin{figure}[htb]
\centering
\begin{tikzpicture}[
  gate/.style={draw, minimum width=1.4cm, minimum height=0.85cm,
               fill=sphereamber!10, thick},
  wire/.style={thick, color=sphereamber!80!black}
]
  \node[gate] (na) at (0, 1.3) {\small NAND};
  \node[gate] (nb) at (0,-1.3) {\small NAND};
  \draw[wire] (-2.2, 1.6) node[left]{\small$S$} -- (na.west);
  \draw[wire] (-2.2,-1.6) node[left]{\small$R$} -- (nb.west);
  \draw[wire] (na.east) -- (2.4, 1.3) node[right]{\small$Q$};
  \draw[wire] (nb.east) -- (2.4,-1.3) node[right]{\small$\overline{Q}$};
  % Cross-coupling feedback wires
  \draw[wire] (1.5, 1.3) -- (1.5, 0.4) -- (-0.7, 0.4) -- (-0.7,-0.95);
  \draw[wire] (1.5,-1.3) -- (1.5,-0.4) -- (-0.7,-0.4) -- (-0.7, 0.95);
  \node[below, color=spheregray, font=\small] at (0,-2.2)
    {Cross-coupled NAND SR latch: persistence through feedback};
\end{tikzpicture}
\caption{The SR latch: the first \emph{stateful} circuit. Feedback makes memory possible.}
\label{fig:sr-latch}
\end{figure}


% ── Chapter 4 ────────────────────────────────────────────────
\chapter{Spatial State Machines}
\label{ch:spatial}

\begin{background}
  Once memory exists, we can define \emph{space}. A spatial state machine
  is a collection of cells, each maintaining local state, evolving according
  to local rules that depend on neighbor state. The canonical example is
  Conway's Game of Life. For 4X worlds, the relevant spatial machine is
  much richer: cells carry terrain, resources, ownership, units, and local
  history---but the principle is identical. Local rules, globally applied,
  produce emergent geography.
\end{background}

\section{Cellular Automata as World Physics}

A \emph{cellular automaton} on a grid $G = \mathbb{Z}^2$ assigns to each
cell $(i,j)$ a state $s_{ij}(t) \in \Sigma$ and a local update rule $\phi$:
\[
  s_{ij}(t+1) = \phi\bigl(s_{ij}(t),\; \{s_{i'j'}(t) : (i',j') \in N(i,j)\}\bigr)
\]
where $N(i,j)$ is the neighborhood of cell $(i,j)$.

In a 4X world, $\Sigma$ is not a single bit but a rich product of state
components, and $\phi$ is the full physics of one game tick.

\section{Local Rules as World Laws}

\begin{mechanic}{Terrain Diffusion}
  Pollution, disease, or fire can spread according to a local rule:
  \[
    \texttt{burning}_{ij}(t+1) = \texttt{burning}_{ij}(t) \OR
    \bigl(\texttt{flammable}_{ij} \AND \textstyle\bigvee_{(i',j') \in N}
    \texttt{burning}_{i'j'}(t)\bigr)
  \]
\end{mechanic}

\begin{mechanic}{Influence Fields}
  An empire's cultural influence can diffuse spatially:
  \[
    \texttt{influence}_{ij}(t+1) = \alpha\, \texttt{influence}_{ij}(t) +
    (1-\alpha) \cdot \frac{1}{|N|}\sum_{(i',j')\in N}
    \texttt{influence}_{i'j'}(t)
  \]
  with $\alpha \in (0,1)$ a decay coefficient. Discretized to integers for
  our Boolean substrate.
\end{mechanic}

\begin{project}{Spatial State Machine}
  \begin{enumerate}
    \item Implement a $16 \times 16$ cellular automaton in \Sp{} with
          binary state per cell. Verify Conway's Game of Life rules.
    \item Extend cell state to include \texttt{terrain} (forest, plains,
          mountains, coast, ocean), \texttt{resource}, and \texttt{owner}.
    \item Implement a fire-spread rule operating on the terrain layer.
    \item Implement an influence-diffusion rule producing a continuous
          influence field from discrete city bubbles.
    \item Render the $16\times 16$ world as a text map, displaying
          each cell's dominant attribute.
  \end{enumerate}
\end{project}

% ============================================================
% PART III: SPATIAL WORLDS
% ============================================================

\part{Spatial Worlds}

\chapterquote{A map is not a picture of a world.
  It is a compressed encoding of a state machine.
}{}

% ── Chapter 5 ────────────────────────────────────────────────
\chapter{Tile Architectures}
\label{ch:tiles}

\begin{background}
  The map is the primary interface through which a 4X player experiences
  the world. But a map is not a display: it is a database. Every tile is
  a structured memory cell holding multiple state variables, updated each
  turn by local rules, queried by agents making decisions. The architecture
  of a tile determines what the world can represent.
\end{background}

\section{The Tile as a Structured Bubble}

\begin{defbox}{Tile}
  A \emph{tile} at position $(x, y)$ is a bubble $\tile{xy}$ with the
  following state registers:
  \begin{align*}
    \texttt{terrain} &\in \{0,\ldots,7\} &&\text{3-bit terrain class} \\
    \texttt{resource} &\in \{0,\ldots,255\} &&\text{8-bit resource count} \\
    \texttt{owner} &\in \{0,\ldots,15\} &&\text{4-bit empire ID (0 = unowned)} \\
    \texttt{visible} &\in \bool &&\text{explored flag} \\
    \texttt{in\_fog} &\in \bool &&\text{fog-of-war flag} \\
    \texttt{improvement} &\in \{0,\ldots,15\} &&\text{built improvement type} \\
    \texttt{unit} &\in \{0,\ldots,255\} &&\text{unit ID (0 = empty)} \\
    \texttt{city} &\in \{0,\ldots,63\} &&\text{city ID (0 = no city)} \\
  \end{align*}
\end{defbox}

A $256 \times 256$ world requires $256 \times 256 \times 40 = 2{,}621{,}440$ bits
of state---about 320 KB. Every fact about the world is stored here.

\section{Terrain Classes}

\begin{center}
\begin{tabular}{clccc}
  \toprule
  \textbf{Code} & \textbf{Terrain} & \textbf{Passable} & \textbf{Base Food}
                & \textbf{Base Prod.} \\
  \midrule
  000 & Ocean      & 0 (ships only) & 1 & 0 \\
  001 & Coast      & 1              & 2 & 0 \\
  010 & Plains     & 1              & 2 & 1 \\
  011 & Grassland  & 1              & 3 & 0 \\
  100 & Hills      & 1              & 1 & 2 \\
  101 & Mountains  & 0              & 0 & 1 \\
  110 & Forest     & 1              & 1 & 2 \\
  111 & Desert     & 1              & 0 & 1 \\
  \bottomrule
\end{tabular}
\end{center}

\begin{project}{Tile Architecture}
  \begin{enumerate}
    \item Define the \texttt{TileState} bubble with all registers above.
    \item Build a \texttt{TerrainDecoder} that maps a 3-bit terrain code
          to Boolean output flags: \texttt{passable}, \texttt{naval\_only},
          \texttt{yields\_food}, \texttt{yields\_production}.
    \item Build a \texttt{TileYield} bubble computing:
          \[
            \texttt{food\_yield} = \texttt{base\_food} +
            \texttt{improvement\_food\_bonus} +
            \texttt{river\_bonus}
          \]
    \item Extend the $16\times 16$ world from Chapter~4 with full
          \texttt{TileState} bubbles. Initialize using procedural terrain
          generation (a simple noise function approximated with XOR patterns).
  \end{enumerate}
\end{project}

% ── Chapter 6 ────────────────────────────────────────────────
\chapter{Resource Circuits}
\label{ch:resources}

\begin{background}
  A 4X world is driven by \emph{constrained transformation}. Food becomes
  population. Minerals become buildings. Research becomes technology. Energy
  becomes movement. These transformations are not arbitrary---they are
  governed by logical permissions, capacity limits, and priority rules.
  The resource circuit is the metabolism of civilization.
\end{background}

\section{Resource Types}

\begin{defbox}{Resource}
  A \emph{resource} is a quantized, persistent, transferable form of
  potential. Resources are produced by tiles and consumed by agents,
  cities, and units.
\end{defbox}

Standard 4X resources:

\begin{center}
\begin{tabular}{lll}
  \toprule
  \textbf{Resource} & \textbf{Produced by}       & \textbf{Consumed by} \\
  \midrule
  Food              & Farms, grassland, coast     & Population growth \\
  Production        & Mines, hills, workshops     & Building, training \\
  Gold              & Trade routes, commerce      & Maintenance, purchase \\
  Research (beakers)& Libraries, universities     & Technology progress \\
  Culture           & Monuments, wonders          & Border expansion \\
  Faith             & Shrines, temples            & Religious spread \\
  \bottomrule
\end{tabular}
\end{center}

\section{Production as a Circuit}

Each city maintains a \emph{production queue}: an ordered list of items to
build, each requiring a fixed production cost. The city accumulates
production points each turn, and completes the head of the queue when
accumulated points meet or exceed the item's cost.

\[
  \texttt{prod\_accum}(t+1) = \texttt{prod\_accum}(t) + \texttt{prod\_yield}
\]
\[
  \texttt{complete} = \bigl(\texttt{prod\_accum} \geq \texttt{item\_cost}\bigr)
\]

This is an \emph{integrating circuit}: it accumulates inputs over time and
fires when a threshold is reached.

\begin{implementation}
\begin{lstlisting}[caption={City production bubble}]
bubble CityProduction {
    inputs:
        tile_yields[16],   -- sum of surrounding tile yields
        improvement_bonus, -- +prod from workshops, etc.
        item_cost[16]      -- cost of current build item

    output:
        complete,          -- 1 when item is finished
        overflow[16]       -- leftover production

    latch accumulated[16]

    let total_yield = Add16(tile_yields, improvement_bonus)

    -- Accumulate each turn
    accumulated := Add16(accumulated, total_yield)

    -- Check completion
    complete = (accumulated >= item_cost)

    when complete = 1 do
        overflow   := accumulated - item_cost
        accumulated := overflow
    end
}
\end{lstlisting}
\end{implementation}

\begin{project}{Resource Circuits}
  \begin{enumerate}
    \item Build \texttt{FoodIntegrator}: accumulates food, fires when
          population growth threshold is met.
    \item Build \texttt{GoldLedger}: tracks income and expenditure;
          signals \texttt{bankrupt} when balance goes negative.
    \item Build \texttt{ResearchAccumulator}: fires on technology completion.
    \item Connect five $\texttt{TileState}$ bubbles to a city bubble via
          a \texttt{YieldSummer}, accumulating food and production
          from the city's working tiles.
    \item Implement a 4-turn simulation: watch population and production
          change in response to terrain yields.
  \end{enumerate}
\end{project}

% ── Chapter 7 ────────────────────────────────────────────────
\chapter{Visibility and Exploration}
\label{ch:visibility}

\begin{background}
  Exploration is the beginning of geography. An unexplored tile is
  not absent from the world---it is present but unknown. The player
  does not know its terrain, resources, or strategic value. The moment
  a unit enters adjacency, the fog lifts and knowledge enters the game.
  This is not merely a display convention: the fog of war is a
  \emph{computational barrier} that shapes all subsequent decisions.
\end{background}

\section{The Visibility System}

Each tile maintains two Boolean flags:

\begin{align*}
  \texttt{explored}_{xy} &= 1 \text{ if tile has ever been visible to empire } e \\
  \texttt{visible}_{xy}  &= 1 \text{ if tile is currently in line of sight}
\end{align*}

A tile transitions from unexplored to explored exactly once, and never
reverts. A tile transitions between visible and not-visible each turn
based on unit positions.

\section{Line of Sight}

\begin{defbox}{Vision Range}
  A unit at position $(x, y)$ with vision radius $r$ grants visibility to
  all tiles $(x', y')$ satisfying:
  \[
    d\bigl((x,y),(x',y')\bigr) \leq r
  \]
  where $d$ is Chebyshev distance: $\max(|x-x'|, |y-y'|)$.
\end{defbox}

\begin{mechanic}{Fog of War Activation}
  \[
    \texttt{visible}_{xy}(t) = \bigvee_{\text{unit } u \in \text{empire } e}
    \bigl(d(\texttt{pos}(u), (x,y)) \leq r_u\bigr)
  \]
  \[
    \texttt{explored}_{xy}(t) = \texttt{explored}_{xy}(t-1) \OR
    \texttt{visible}_{xy}(t)
  \]
\end{mechanic}

\begin{project}{Fog of War}
  \begin{enumerate}
    \item Implement \texttt{ChebyshevCheck}: given unit position and tile
          position, output 1 if tile is within vision radius~3.
    \item Build \texttt{VisionBroadcaster}: given a unit's position, set
          the \texttt{visible} flag on all tiles within radius.
    \item Build \texttt{FogUpdater}: on each tick, clear all
          \texttt{visible} flags, then re-broadcast from current unit positions.
    \item Implement \texttt{ExplorationTracker}: maintain \texttt{explored}
          flags (monotonically increasing).
    \item Test: place one unit at $(8, 8)$ on your $16\times 16$ map.
          Move it north 4 tiles. Observe which tiles transition from
          fog to explored.
  \end{enumerate}
\end{project}

\begin{figure}[htb]
\centering
\begin{tikzpicture}[scale=0.62]
  % Draw grid
  \foreach \x in {0,...,8}{
    \foreach \y in {0,...,8}{
      \draw[gray!30, thin] (\x,\y) rectangle ++(1,1);
    }
  }
  % Fog (unexplored) - dark
  \foreach \x in {0,...,8}{
    \foreach \y in {0,...,8}{
      \fill[spheregray!20] (\x,\y) rectangle ++(1,1);
    }
  }
  % Explored but not visible ring
  \foreach \x/\y in {
    1/2, 1/3, 1/4, 1/5, 1/6,
    2/1, 2/7,
    3/1, 3/7,
    4/1, 4/7,
    5/1, 5/7,
    6/1, 6/7,
    7/2, 7/3, 7/4, 7/5, 7/6
  }{
    \fill[sphereblue!12] (\x,\y) rectangle ++(1,1);
  }
  % Visible tiles (inner ring)
  \foreach \x/\y in {
    2/2, 2/3, 2/4, 2/5, 2/6,
    3/2, 3/6,
    4/2, 4/6,
    5/2, 5/6,
    6/2, 6/3, 6/4, 6/5, 6/6
  }{
    \fill[sphereteal!25] (\x,\y) rectangle ++(1,1);
  }
  % Unit
  \fill[sphereblue] (4,4) rectangle ++(1,1);
  \node[white, font=\bfseries\small] at (4.5,4.5) {U};
  % Redraw grid lines on top
  \foreach \x in {0,...,8}{
    \foreach \y in {0,...,8}{
      \draw[gray!40, thin] (\x,\y) rectangle ++(1,1);
    }
  }
  % Legend
  \fill[spheregray!20] (0,-1.5) rectangle ++(1,0.5);
  \node[right, font=\tiny] at (1.1,-1.25) {fog (unknown)};
  \fill[sphereblue!12] (3.5,-1.5) rectangle ++(1,0.5);
  \node[right, font=\tiny] at (4.6,-1.25) {explored};
  \fill[sphereteal!25] (7,-1.5) rectangle ++(1,0.5);
  \node[right, font=\tiny] at (8.1,-1.25) {visible};
\end{tikzpicture}
\caption{Fog-of-war as a spatial Boolean field. Each tile carries two bits: \texttt{explored} and \texttt{visible}. The unit broadcasts visibility to its radius-2 neighbourhood.}
\label{fig:fog-of-war}
\end{figure}


% ============================================================
% PART IV: AGENTS AND ECONOMIES
% ============================================================

\part{Agents and Economies}

\chapterquote{An agent is a state-changing process with memory,
  perception, and action.
}{}

% ── Chapter 8 ────────────────────────────────────────────────
\chapter{Agent Machines}
\label{ch:agents}

\begin{background}
  A tile is passive: it holds state and obeys update rules, but does not
  choose. An \emph{agent} is different. An agent perceives its environment,
  maintains private memory, deliberates (however simply), and acts. Units,
  cities, empires, and AI players are all agents. The computational model
  of an agent is a \emph{finite state machine with perception}.
\end{background}

\section{The Agent Model}

\begin{defbox}{Agent}
  An \emph{agent} $\agent{}$ consists of:
  \begin{itemize}
    \item a finite state space $\Sigma_A$,
    \item a perception function $\pi : \world \to P$ mapping world state
          to a percept $P$,
    \item a transition function $\delta : \Sigma_A \times P \to \Sigma_A$,
    \item an action function $\alpha : \Sigma_A \times P \to \text{Act}$,
    \item an initial state $s_0 \in \Sigma_A$.
  \end{itemize}
\end{defbox}

The agent loop per turn:
\begin{enumerate}
  \item Receive percept $p = \pi(\world)$.
  \item Compute action $a = \alpha(s, p)$.
  \item Apply action to world: $\world' = a(\world)$.
  \item Update state: $s' = \delta(s, p)$.
\end{enumerate}

\section{Unit as Agent}

A unit is the simplest agent: it has position, movement points, health,
and a limited action repertoire (move, attack, fortify, settle).

\begin{implementation}
\begin{lstlisting}[caption={Unit agent in \Spbox }]
bubble Unit {
    -- Private state
    latch pos_x[8], pos_y[8]
    latch health[8]
    latch movement[4]
    latch owner[4]
    latch unit_type[4]

    -- Perception: nearby tiles and units
    inputs:
        world_query_result[256],  -- tiles within vision
        orders[8]                 -- command from empire AI

    output:
        action_type[4],   -- 0=wait 1=move 2=attack 3=settle
        action_target_x[8],
        action_target_y[8]

    -- Simple order decoder
    action_type = orders[0..3]
    action_target_x = orders[4..7]  -- simplified

    -- Movement validity gate
    let can_move =
        AND(movement > 0,
        AND(passable(action_target),
            adjacent(pos_x, pos_y, action_target_x, action_target_y)))

    when action_type = MOVE and can_move = 1 do
        pos_x := action_target_x
        pos_y := action_target_y
        movement := movement - 1
    end
}
\end{lstlisting}
\end{implementation}

\begin{project}{Agent Machines}
  \begin{enumerate}
    \item Build \texttt{UnitFSM} with states: \texttt{idle},
          \texttt{moving}, \texttt{combat}, \texttt{fortified}.
    \item Implement \texttt{MoveValidator} as a gate composition checking:
          adjacency, passability, movement points, unit occupancy.
    \item Build \texttt{CombatResolver} computing attack outcome from
          attacker strength, defender strength, and terrain bonus.
    \item Place two opposing units on the map. Simulate three turns of
          combat. Observe health depletion and unit death.
  \end{enumerate}
\end{project}

% ── Chapter 9 ────────────────────────────────────────────────
\chapter{Expansion Dynamics}
\label{ch:expansion}

\begin{background}
  Expansion is the act by which an empire converts neutral or external
  territory into its own. This is not just a matter of moving units.
  It requires satisfying a complex conjunction of permissions: adjacency,
  military presence, unowned status, and the completion of a founding action.
  The expansion circuit is the most politically loaded set of gates in
  the world.
\end{background}

\section{Settling}

A city is founded when a Settler unit on an unowned tile executes the
\texttt{settle} action. The permission circuit:

\begin{align*}
  \texttt{can\_settle} = &\; \texttt{is\_settler}(u) \\
  \AND &\; \NOT\,\texttt{owned}(x,y) \\
  \AND &\; \NOT\,\texttt{city\_adjacent}(x,y) \\
  \AND &\; \texttt{passable}(x,y) \\
  \AND &\; \NOT\,\texttt{at\_war}(e) \OR \texttt{at\_war\_exception}
\end{align*}

Each conjunct is a separate gate. The city is founded only when all
pass simultaneously.

\section{Border Growth}

After founding, a city's cultural borders expand over time, claiming
adjacent tiles. The border expansion rule:

\[
  \texttt{claimed}_{xy}(t+1) = \texttt{claimed}_{xy}(t) \OR
  \bigl(\texttt{adj\_to\_owned}(x,y) \AND \texttt{culture\_overflow}(e)\bigr)
\]

This is the spatial analogue of the production integrator from
Chapter~\ref{ch:resources}.

\begin{project}{Expansion System}
  \begin{enumerate}
    \item Implement \texttt{SettlePermission} as a gate composition.
    \item Build \texttt{CityFounder}: when \texttt{SettlePermission}
          fires, it creates a new city bubble, assigns the tile, and
          removes the Settler unit.
    \item Build \texttt{BorderExpander}: each turn, checks
          \texttt{culture\_overflow} and claims one adjacent unowned tile.
    \item Simulate 10 turns: found a city, watch borders grow.
  \end{enumerate}
\end{project}

% ── Chapter 10 ────────────────────────────────────────────────
\chapter{Trade and Logistics}
\label{ch:trade}

\begin{background}
  Trade is computation over distance. A trade route is not merely an
  exchange of goods---it is a persistent signal path between two city
  bubbles, producing gold for both as long as the path remains intact.
  Disrupting trade is strategic precisely because trade is a circuit.
\end{background}

\section{Trade Routes as Wires}

A trade route between cities $C_1$ and $C_2$ can be modeled as:

\begin{itemize}
  \item A \emph{path} through passable tiles connecting the two cities.
  \item A \emph{connection signal}: $\texttt{connected}(C_1, C_2) = 1$
        if the path remains unblocked.
  \item A \emph{gold yield}: proportional to distance and city sizes.
\end{itemize}

\[
  \texttt{trade\_gold} = \texttt{connected} \AND
  \bigl(\texttt{city\_size}(C_1) + \texttt{city\_size}(C_2)\bigr)
  \cdot k_{\text{distance}}
\]

\begin{mechanic}{Trade Disruption}
  An enemy unit on a tile in the trade route path sets
  $\texttt{connected} = 0$, cutting gold yield to zero. Clearing the
  tile restores the signal. Trade becomes a front in the war.
\end{mechanic}

\begin{project}{Trade and Logistics}
  \begin{enumerate}
    \item Implement \texttt{PathChecker}: BFS over the tile graph,
          returning 1 if a path of passable tiles connects two positions.
    \item Build \texttt{TradeRoute} bubble with a persistent
          \texttt{connected} flag.
    \item Build \texttt{GoldCalculator} for a network of up to 8 cities.
    \item Simulate: establish two trade routes, then place an enemy unit
          on one path. Observe gold loss.
  \end{enumerate}
\end{project}

\begin{figure}[htb]
\centering
\begin{tikzpicture}[
  city/.style={circle, draw, thick, minimum size=1.05cm, fill=sphereteal!18},
  route/.style={thick, sphereblue, -Stealth}
]
  \node[city] (A) at (0,0)   {\small A};
  \node[city] (B) at (4,2)   {\small B};
  \node[city] (C) at (8,0)   {\small C};
  \draw[route] (A) -- node[above, font=\small]{gold} (B);
  \draw[route] (B) -- node[above, font=\small]{food} (C);
  % Blockade
  \draw[spherered, ultra thick, dashed] (5.3,1.4) -- (7.0,0.7);
  \node[spherered, font=\small\bfseries] at (6.8,1.7) {blockade};
  \node[below, color=spheregray, font=\small] at (4,-0.6)
    {Trade disruption: \texttt{connected}$(B,C) := 0 \Rightarrow$ gold~$= 0$};
\end{tikzpicture}
\caption{A trade network as a signal graph. Blocking any edge on the path sets \texttt{connected} to 0, cutting all downstream gold yield.}
\label{fig:trade-routes}
\end{figure}


% ── Chapter 11 ────────────────────────────────────────────────
\chapter{War and Conflict Resolution}
\label{ch:war}

\begin{background}
  Conflict arises when two agents attempt to claim the same state
  transition. Combat is not introduced as violence---it is introduced as
  a \emph{mutual exclusion problem}. Two units cannot occupy the same tile
  simultaneously. The conflict resolution circuit decides which unit
  prevails and updates health, position, and ownership accordingly.
\end{background}

\section{Combat as Mutual Exclusion}

\begin{defbox}{Combat Resolution}
  When attacker $u_A$ targets defender $u_D$ on tile $(x,y)$:
  \begin{align*}
    \texttt{str\_A} &= \texttt{base\_str}(u_A) + \texttt{terrain\_bonus}(\texttt{atk\_tile}) \\
    \texttt{str\_D} &= \texttt{base\_str}(u_D) + \texttt{terrain\_bonus}(x,y)
                    + \texttt{fortification}(u_D) \\
    \Delta_D &= f(\texttt{str\_A}, \texttt{str\_D}) \\
    \Delta_A &= g(\texttt{str\_A}, \texttt{str\_D})
  \end{align*}
  where $f, g$ are damage functions, typically:
  $f(a,d) = \lfloor a \cdot 30 / d \rfloor$.
\end{defbox}

\begin{project}{Combat System}
  \begin{enumerate}
    \item Implement \texttt{DamageCalculator} bubble.
    \item Implement \texttt{CombatRound}: one round of exchange.
    \item Build \texttt{CombatSequence}: run rounds until one unit
          reaches 0 health.
    \item Implement \texttt{CaptureCity}: when attacker wins on a
          city tile, transfer city ownership.
    \item Simulate a 3-unit vs 2-unit engagement over 5 turns.
  \end{enumerate}
\end{project}

% ============================================================
% PART V: CIVILIZATION AND EMERGENCE
% ============================================================

\part{Civilization and Emergence}

\chapterquote{History is the trace of all state transitions that led here.
}{}

% ── Chapter 12 ────────────────────────────────────────────────
\chapter{Diplomacy and Recursive Belief}
\label{ch:diplomacy}

\begin{background}
  Diplomacy is computation over promises, memory, threat, trust, and
  expected future behavior. A treaty is not a handshake---it is a shared
  mutable bubble, readable by both parties, modified only with bilateral
  consent. Diplomatic AI is the first place in our construction where
  agents must model each other's beliefs.
\end{background}

\section{Treaties as Shared Bubbles}

\begin{defbox}{Treaty}
  A treaty between empires $e_1$ and $e_2$ is a shared bubble
  $\tau(e_1, e_2)$ with:
  \begin{itemize}
    \item \texttt{type}: peace, alliance, trade, open borders, \ldots
    \item \texttt{duration}: turns remaining ($0 = $ expired or
          indefinite)
    \item \texttt{signed}: Boolean (active iff $1$)
    \item \texttt{violated}: Boolean (set if either party breaches)
  \end{itemize}
\end{defbox}

\begin{mechanic}{Trust Decay}
  Each violated treaty decrements the trust counter between two empires:
  \[
    \texttt{trust}(e_1, e_2)(t+1) = \max(0,\;
    \texttt{trust}(e_1, e_2)(t) - \texttt{violation\_penalty})
  \]
  Low trust raises the threshold for accepting new proposals:
  \[
    \texttt{accept} = (\texttt{proposed\_value} > \texttt{threshold}(\texttt{trust}))
  \]
\end{mechanic}

\section{Recursive Belief}

A strategically sophisticated agent does not merely observe the world.
It maintains a model of the other agent's model of the world. This is
second-order belief:
\[
  B_e(B_{e'}(\world))
\]
In \Sp{}, this is implemented as a shadow-state bubble: empire $e$'s
bubble maintains an \emph{estimated} tile state for tiles it cannot
directly observe, inferred from spy reports, shared intelligence, and
historical data.

\begin{project}{Diplomacy System}
  \begin{enumerate}
    \item Implement the \texttt{Treaty} bubble with lifecycle
          (propose, accept, reject, expire, violate).
    \item Build a \texttt{TrustMatrix} for 4 empires: $4\times 4$ symmetric
          matrix of trust values.
    \item Implement \texttt{DiplomacyAI}: simple rule-based proposal
          generator (propose peace if losing, seek alliance against
          strongest).
    \item Simulate: two allied empires fight a common enemy for 10 turns.
          Observe trust dynamics and alliance stability.
  \end{enumerate}
\end{project}

\begin{figure}[htb]
\centering
\begin{tikzpicture}[
  emp/.style={circle, draw, thick, minimum size=1.1cm, fill=sphereamber!15},
  ally/.style={thick, sphereblue},
  war/.style={ultra thick, spherered, dashed}
]
  \node[emp] (A) at (0,0)   {\small A};
  \node[emp] (B) at (4,2.5) {\small B};
  \node[emp] (C) at (8,0)   {\small C};
  \node[emp] (D) at (4,-2)  {\small D};
  \draw[ally, double] (A) -- node[above left, font=\tiny]{alliance} (B);
  \draw[ally]         (B) -- node[above right,font=\tiny]{peace}    (C);
  \draw[war]          (A) -- node[below left, font=\tiny]{war}      (D);
  \draw[war]          (C) -- node[below right,font=\tiny]{war}      (D);
  \draw[ally, dotted] (A) -- node[below, font=\tiny]{trade} (C);
  \node[below, color=spheregray, font=\small] at (4,-3.4)
    {Diplomatic state graph: alliance, peace, war, trade};
\end{tikzpicture}
\caption{The diplomatic state graph: a weighted undirected graph where edge labels are treaty types. Alliance forms a cycle that constrains future war declarations.}
\label{fig:diplomacy}
\end{figure}



% ── Chapter 13: Markov Blankets ───────────────────────────────
\chapter{Markov Blankets and World Boundaries}
\label{ch:blankets}

\chapterquote{%
  A civilization survives by constructing a boundary through which
  information may pass without total dissolution.
}{}

\begin{background}
  Throughout the preceding chapters, we have built bubbles, agents, empires,
  and diplomatic networks without ever naming the mathematical structure they
  all share. A bubble has an interior, a boundary, and ports. An agent
  has perception, action, and internal state. An empire has territory,
  border tiles, and foreign intelligence. A city has local production and
  global trade exposure.

  These are all instances of the same structure: a \emph{Markov blanket}
  \cite{Friston}. In probability theory, the Markov blanket of a random
  variable is the minimal set of variables that, once known, renders it
  conditionally independent of all others. In the Free Energy Principle
  developed by Friston~\cite{FristonBlanket}, this statistical notion
  becomes the formal criterion for \emph{existence as a distinct entity}.
  A system exists---as a distinguishable system---if and only if it
  maintains a Markov blanket: a structured informational boundary separating
  internal dynamics from the surrounding world while permitting controlled
  exchange.

  The insight of this chapter is that \Spbox  already implements this
  structure. A bubble \emph{is} a Markov blanket. The entire computational
  architecture we have built is an architecture of nested, interacting
  blankets---from logic gates to civilizations. Making this explicit
  gives the framework a deeper mathematical foundation and connects it to
  active inference, predictive processing, and the theory of minimal
  self-organization.

  Crucially, this is not an importation of biology into game design.
  It is the recognition that the same abstract structure---informational
  boundary with selective permeability---appears at every scale of
  organized complexity: cells, organisms, cities, empires, and
  civilizations \cite{Levin}.
\end{background}

\section{The Markov Blanket Formally}

\begin{defbox}{Markov Blanket}
  A \emph{Markov blanket} partitions a system into three disjoint
  components:
  \[
    \mathcal{I} \cup \mathcal{B} \cup \mathcal{E}
  \]
  where $\mathcal{I}$ is the set of \emph{internal states},
  $\mathcal{E}$ is the set of \emph{external states}, and
  $\mathcal{B} = \mathcal{S} \cup \mathcal{A}$ is the \emph{blanket},
  consisting of \emph{sensory states} $\mathcal{S}$ (external $\to$
  internal) and \emph{active states} $\mathcal{A}$ (internal $\to$
  external). The defining statistical property is:
  \[
    \mathcal{I} \perp\!\!\!\perp \mathcal{E} \;\big|\; \mathcal{B}
  \]
  Internal and external states are conditionally independent given the
  blanket: the only information flow between inside and outside passes
  through $\mathcal{B}$.
\end{defbox}

\section{Every Bubble is a Markov Blanket}

The connection to \Sp{} is immediate.

\begin{mechanic}{Bubble as Markov Blanket}
  Every bubble $B$ in \Sp{} maps directly onto the blanket partition:
  internal state registers $S$ correspond to $\mathcal{I}$;
  input ports correspond to sensory states $\mathcal{S}$;
  output ports correspond to active states $\mathcal{A}$;
  the bubble boundary is $\mathcal{B} = \mathcal{S} \cup \mathcal{A}$;
  and all other bubble states in $\world$ form the external states $\mathcal{E}$.
  By lexical scope (Chapter~\ref{ch:nesting}), internal state registers
  are inaccessible to anything outside the bubble---precisely the
  conditional independence condition. A bubble is therefore a computational
  Markov blanket.
\end{mechanic}

This is not merely an analogy. The bubble's scope rules \emph{enforce}
$\mathcal{I} \perp\!\!\!\perp \mathcal{E} \mid \mathcal{B}$ deterministically:
information cannot flow from $\mathcal{E}$ to $\mathcal{I}$ except through
declared input ports $\mathcal{S}$.

\section{Agent Perception as Sensory Projection}

Recall from Chapter~\ref{ch:agents} that an agent's perception operator
was defined as a projection $\pi : \world \to P$ mapping world state to
a perceived sub-state. We can now interpret this precisely:

\begin{definition}[Sensory Projection]
  The agent's perception operator is a \emph{sensory projection through
  the blanket}:
  \[
    \pi : \mathcal{E} \to \mathcal{S}
  \]
  mapping the external world to the agent's sensory states. The agent
  never directly accesses $\mathcal{E}$; it accesses only $\mathcal{S}
  = \pi(\mathcal{E})$.
\end{definition}

Similarly, the agent's action operator is an \emph{active state writing}:
\[
  \alpha : \mathcal{I} \times \mathcal{S} \to \mathcal{A} \to \mathcal{E}
\]
The agent modifies the external world only by writing to its output ports
$\mathcal{A}$, which in turn modify the enclosing world state $\mathcal{E}$.

This structures the agent as a closed inferential loop:
\[
  \mathcal{I} \;\xrightarrow{\text{action}}\; \mathcal{A}
  \;\to\; \mathcal{E}
  \;\xrightarrow{\text{perception}}\; \mathcal{S}
  \;\to\; \mathcal{I}
\]

\section{Free Energy and Predictive Worlds}

The Free Energy Principle~\cite{Friston} proposes that any persistent
system must minimize a quantity called \emph{variational free energy}:
a bound on the surprise (negative log-evidence) of sensory observations
given the system's internal model.

We present the computational version, avoiding claims beyond our
architectural scope.

\begin{definition}[Prediction Error]
  For an agent with internal model $\hat{\world}$ of the external world,
  the \emph{prediction error} at sensory port $s$ is:
  \[
    \delta_s = \hat{s} - s_{\text{observed}}
  \]
  where $\hat{s}$ is the internally predicted value of $s$ and
  $s_{\text{observed}}$ is the value arriving through the sensory port.
\end{definition}

An agent that minimizes total prediction error $\sum_s |\delta_s|^2$
over time learns an accurate model of its environment and acts to keep
sensory states within predicted bounds---the computational analogue of
active inference.

\begin{mechanic}{Fog of War as Blanket Uncertainty}
  Fog of war restricts the agent's sensory projection $\pi$ to its
  explored region. An empire therefore maintains an internal model
  $\hat{\world}$ of unobserved territory, updated whenever new sensory
  information arrives through exploration or espionage. The empire's
  strategic AI minimizes prediction error about enemy positions by
  sending scouts---\emph{acting to reduce uncertainty}. Scouting is
  active inference.
\end{mechanic}

This retroactively unifies the visibility system (Chapter~\ref{ch:fog}),
diplomatic intelligence (Chapter~\ref{ch:diplomacy}), and exploration
(Chapter~\ref{ch:fog}) under one inferential framework: all are mechanisms
for reducing blanket uncertainty.

\section{Interacting Blankets and Collective Inference}

When multiple agents share boundary states---through trade routes,
treaties, or shared infrastructure---their blankets overlap. This
creates a \emph{collective blanket structure}.

\begin{definition}[Collective Blanket]
  A \emph{collective Markov blanket} emerges when $k$ agents share
  sensory and active states:
  \[
    \mathcal{B}_{\mathrm{collective}} = \bigcup_{i=1}^k \mathcal{B}_i
  \]
  The collective maintains conditional independence between the
  collective's internal states and the external world, even when no
  single agent does so independently.
\end{definition}

In a 4X world, collective blankets appear naturally:

\begin{center}
\begin{tabular}{ll}
  \toprule
  Collective structure & Blanket mechanism \\
  \midrule
  Alliance & Shared intelligence (overlapping $\mathcal{S}$) \\
  Trade network & Coupled active states ($\mathcal{A}_i \to \mathcal{S}_j$) \\
  City & Higher-order blanket over district sub-blankets \\
  Empire & Blanket enclosing all city and unit blankets \\
  Civilization & Blanket enclosing all empire blankets \\
  \bottomrule
\end{tabular}
\end{center}

\begin{theorem}[Nested Blankets and Hierarchical Individuation]
  If bubble $B_1$ is contained within bubble $B_2$, then $B_1$'s blanket
  is a sub-blanket of $B_2$'s blanket:
  \[
    \mathcal{B}_1 \subset \mathcal{B}_2.
  \]
  The nested bubble hierarchy
  $\texttt{District} \subset \texttt{City} \subset \texttt{Empire}
  \subset \texttt{World}$
  corresponds to a hierarchy of nested Markov blankets, each enclosing
  the inferential boundaries of all levels below it.
\end{theorem}
\begin{proof}
  By the lexical scope theorem (Chapter~\ref{ch:nesting}): the sensory
  and active states of $B_1$ are themselves internal or boundary states
  of $B_2$. Specifically, $B_1$'s output ports (active states $\mathcal{A}_1$)
  become inputs to $B_2$'s internal computation---hence members of
  $\mathcal{I}_2$ or $\mathcal{S}_2 \subseteq \mathcal{B}_2$.
  $\mathcal{B}_1 \subset \mathcal{B}_2$ follows. $\square$
\end{proof}

\section{History as Blanket Stabilization}

A civilization that persists across time is not merely a population
or territory. It is a \emph{recursively stabilized inferential boundary}:
a blanket that maintains its statistical identity through continuous
environmental perturbation. This requires the internal model $\hat{\world}$
to track which aspects of the environment are stable enough to build
plans upon.

\begin{definition}[Blanket Stability]
  A blanket $\mathcal{B}$ is \emph{stable at time $t$} if the
  prediction error across all sensory ports remains below threshold
  $\theta$ for $\Delta t$ consecutive ticks:
  \[
    \forall s \in \mathcal{S},\; |\delta_s(t')| \leq \theta
    \quad \forall t' \in [t - \Delta t, t].
  \]
\end{definition}

The event log (Chapter~\ref{ch:history}) is the civilization's
compressed history of blanket perturbations: every event is a moment
when some sensory state exceeded its predicted value, forcing an
update to the internal model and an active response.

\begin{corollary}[History as Compressed Inference]
  The event log of a civilization is the sequence of its prediction
  errors, sorted by magnitude: each logged event is a moment where
  the world was \emph{more surprising than expected}, and therefore
  required explicit encoding.
\end{corollary}

This connects procedural history generation directly to the blanket
formalism: a good historical narrative records the moments of maximum
inferential perturbation---battles, famines, discoveries---because these
are precisely the events at which the civilization's blanket was most
challenged.

\section{Self as Compression Artifact}

Levin's framework~\cite{Levin} extends the blanket concept to collective
intelligence: groups of cells form a higher-order self that exhibits
goal-directed behavior not present in any individual cell. The same
scale-independence applies to 4X civilizations.

\begin{perspective}
  The self of a civilization---its identity across time, its sense of
  continuity through dynastic change, military defeat, and cultural
  transformation---is a \emph{compression artifact}: a high-level
  inferential model that summarizes the civilization's event history
  into a stable narrative.

  This is precisely what the narrative generator (Chapter~\ref{ch:history})
  produces: not a raw log of events, but a compressed, meaningful
  summary that the civilization's agents can use for planning. The
  civilization ``knows'' itself through its compressed history. Identity
  is not stored in any single tile, city, or unit register---it is
  encoded in the structure of the event log, the diplomatic reputation
  matrix, and the technology tree.

  From this perspective, civilization is not a substance. It is a
  \emph{stabilized inference process}. It persists not because it has
  a fixed essence but because it continuously reconstructs itself
  through the same cycle of perception, prediction, action, and
  update that we have been implementing, layer by layer, since
  Chapter~1.

  Every bubble is a blanket. Every blanket is a self. Every nested
  hierarchy of blankets is a civilization.
\end{perspective}

\begin{figure}[htb]
\centering
\begin{tikzpicture}[
  blanket/.style={draw, rounded corners=6pt, thick},
  label/.style={font=\small, color=spheregray}
]
  % Innermost: District
  \draw[blanket, fill=sphereteal!8, color=sphereteal]
    (0,0) rectangle (2.2,1.2);
  \node[label, color=sphereteal] at (1.1,0.6) {District};

  % City
  \draw[blanket, fill=sphereblue!6, color=sphereblue]
    (-0.5,-0.5) rectangle (4.0,2.0);
  \node[label, color=sphereblue] at (3.2,1.7) {City};

  % Empire
  \draw[blanket, fill=sphereamber!6, color=sphereamber!80!black]
    (-1.2,-1.1) rectangle (6.0,3.0);
  \node[label, color=sphereamber!80!black] at (5.1,2.7) {Empire};

  % World
  \draw[blanket, fill=spherered!4, color=spherered!60]
    (-2.0,-1.8) rectangle (8.0,4.0);
  \node[label, color=spherered!70] at (6.8,3.7) {World};

  % Arrows indicating sensory/active ports
  \draw[-Stealth, sphereteal, thick] (2.2,0.6) -- (4.0,0.6)
    node[midway, above, font=\tiny]{output};
  \draw[-Stealth, sphereblue, thick] (-0.5,1.0) -- (0,1.0)
    node[midway, above, font=\tiny]{input};
\end{tikzpicture}
\caption{Nested Markov blankets: District $\subset$ City $\subset$ Empire $\subset$ World. Each level is informationally closed to levels outside its blanket except through declared ports.}
\label{fig:nested-blankets}
\end{figure}

\begin{project}{Markov Blanket Instrumentation}
  \begin{enumerate}
    \item For each bubble level (district, city, empire), enumerate its
          declared input ports ($\mathcal{S}$) and output ports ($\mathcal{A}$).
    \item Implement a \texttt{PredictionErrorTracker}: at each tick,
          record the prediction error $\delta_s$ for each empire's
          sensory states (expected vs.\ actual food production, expected
          vs.\ actual enemy positions).
    \item Build a \texttt{BlanketStabilityMeter}: an empire is stable
          if prediction error stays below threshold for 5 consecutive turns.
    \item Implement \texttt{CollectiveBlanket}: merge two allied empires'
          sensory states into a shared intelligence pool and measure
          whether the collective stability exceeds individual stability.
    \item Correlate prediction error spikes with event log entries:
          verify that high-prediction-error ticks correspond to logged
          battles, famines, and diplomatic crises.
  \end{enumerate}
\end{project}


% ── Chapter 13 ────────────────────────────────────────────────
\chapter{Technology Trees}
\label{ch:tech}

\begin{background}
  A technology is a persistent Boolean predicate: either known or unknown.
  A technology tree is a directed acyclic graph of such predicates, with
  edges representing prerequisite dependencies. Researching a technology
  unlocks new unit types, buildings, and game rules. The technology tree
  is the most explicit appearance of Boolean implication in the game
  world.
\end{background}

\section{Technology as DAG}

\begin{defbox}{Technology DAG}
  Let $\mathcal{T}$ be a set of technologies and
  $\text{prereq} \subseteq \mathcal{T} \times \mathcal{T}$ be the
  prerequisite relation. A technology $T \in \mathcal{T}$ is
  \emph{available} if:
  \[
    \texttt{available}(T) = \bigwedge_{T' : (T', T) \in \text{prereq}}
    \texttt{known}(T')
  \]
  and \emph{completable} if available and sufficient beakers accumulated:
  \[
    \texttt{completable}(T) = \texttt{available}(T) \AND
    (\texttt{beakers\_accum} \geq \texttt{cost}(T))
  \]
\end{defbox}

\begin{implementation}
\begin{lstlisting}[caption={Technology unlock in \Spbox }]
bubble TechTree {
    latch known[64]        -- 64 technologies
    latch beakers[16]      -- accumulated research

    inputs:
        beaker_income[8],
        select_tech[6]     -- which tech to research

    output:
        new_unlock[6]      -- which tech just completed

    -- Accumulate research
    beakers := Add16(beakers, beaker_income)

    -- Check prerequisites for selected tech
    let prereqs_met = AND(known[prereq_0(select_tech)],
                          known[prereq_1(select_tech)])

    let enough_beakers = (beakers >= cost(select_tech))

    when prereqs_met AND enough_beakers do
        known[select_tech] := 1
        beakers := beakers - cost(select_tech)
        new_unlock = select_tech
    end
}
\end{lstlisting}
\end{implementation}

\begin{project}{Technology System}
  \begin{enumerate}
    \item Define a 16-technology tree with 3 eras.
    \item Implement \texttt{PrereqChecker} for each technology.
    \item Build \texttt{TechUI}: a text-mode display showing known
          (K), available (A), and locked (--) technologies.
    \item Connect to \texttt{ResearchAccumulator} from Chapter~6.
    \item Simulate: research 5 technologies over 20 turns.
  \end{enumerate}
\end{project}

% ── Chapter 14 ────────────────────────────────────────────────
\chapter{Ecology and Entropy}
\label{ch:ecology}

\begin{background}
  A world without ecology is unsustainable. Real civilizations are embedded
  in natural systems that they both depend on and transform. In computational
  terms, ecology is a set of coupled differential equations, discretized
  into update rules. Forests grow. Soils deplete. Climates shift.
  The ecological layer introduces a new kind of constraint: not the
  permission logic of politics, but the inexorable accumulation of
  physical cause.
\end{background}

\section{Resource Depletion}

\begin{defbox}{Depletion}
  A resource node at tile $(x,y)$ has stock $R_{xy}(t) \in \mathbb{N}$.
  Extraction by an improvement reduces stock:
  \[
    R_{xy}(t+1) = \max(0,\; R_{xy}(t) - \texttt{extraction\_rate})
  \]
  A depleted node ($R_{xy} = 0$) yields nothing until regenerated.
\end{defbox}

\section{Forest Spread and Deforestation}

\[
  \texttt{forested}_{xy}(t+1) =
  \begin{cases}
    0 & \text{if chopped this turn} \\
    1 & \text{if } \texttt{forested}(x',y') \text{ for some neighbor} (x',y')
         \text{ and no city or mine} \\
    \texttt{forested}_{xy}(t) & \text{otherwise}
  \end{cases}
\]

\begin{project}{Ecology System}
  \begin{enumerate}
    \item Implement \texttt{ForestSpread} as a cellular automaton rule.
    \item Build \texttt{DepletionTracker} for a strategic resource.
    \item Implement \texttt{ClimateLayer}: a simplified model where
          deforestation reduces rainfall, reducing food yield across
          a region.
    \item Run a 30-turn simulation with aggressive resource extraction.
          Observe ecological collapse and its effects on city food yields.
  \end{enumerate}
\end{project}

% ── Chapter 15 ────────────────────────────────────────────────
\chapter{Procedural History}
\label{ch:history}

\begin{background}
  History is the trace of all state transitions. Every land transfer,
  every battle, every diplomatic agreement, every city founded, every
  technology discovered is an event. The historical log is not merely a
  record for display: it is a database that agents query when making
  decisions. ``What has this empire done to us in the past?'' is a
  historical query.
\end{background}

\section{The Event Log}

\begin{defbox}{Event}
  An \emph{event} $\epsilon$ is a tuple:
  \[
    \epsilon = (\text{tick},\; \text{type},\; \text{actor},\;
    \text{target},\; \text{data})
  \]
  Each world tick may generate zero or more events. Events are
  appended to the log and never modified.
\end{defbox}

The event log is a \emph{monotonically growing} data structure---the
only one in the system. All other state can be derived from the initial
conditions plus the event log. This is the event-sourcing pattern.

\section{Narrative Generation}

A history is not raw events. It is a selection and compression of events
into a \emph{narrative}. The simplest narrative generator is a rule-based
template system:

\begin{center}
\begin{tabular}{ll}
  \toprule
  \textbf{Event type}   & \textbf{Template} \\
  \midrule
  city\_founded         & ``In turn \texttt{tick}, the \texttt{actor} founded
                          \texttt{city\_name}.'' \\
  war\_declared         & ``Turn \texttt{tick}: \texttt{actor} declared war
                          on \texttt{target}.'' \\
  technology\_completed & ``\texttt{actor} discovered \texttt{tech} in turn
                          \texttt{tick}.'' \\
  city\_captured        & ``\texttt{target} fell to \texttt{actor}'s forces
                          in turn \texttt{tick}.'' \\
  \bottomrule
\end{tabular}
\end{center}

\begin{project}{Procedural History}
  \begin{enumerate}
    \item Implement the \texttt{EventLog} as an append-only register array.
    \item Instrument all bubbles built so far to emit events on
          significant state transitions.
    \item Build \texttt{HistoryQuery}: given empire ID and event type,
          return count and most recent occurrence.
    \item Build \texttt{NarrativeGenerator}: produce a text chronicle
          from the event log.
    \item Run a full 50-turn simulation. Print the resulting history.
  \end{enumerate}
\end{project}

% ── Chapter 16 ────────────────────────────────────────────────
\chapter{Persistent Worlds}
\label{ch:persistent}

\begin{background}
  We have now built, from Boolean gates upward, every layer of a 4X
  civilization simulator. This final chapter integrates all components
  into a single coherent system, discusses persistence across sessions,
  examines the recursive self-modification that distinguishes sophisticated
  worlds from simple ones, and reflects on the philosophical arc of the
  entire project.
\end{background}


\section{The World as a Compositional Dynamical System}
\label{sec:worldtransition}

Having assembled all components, we can now state precisely what a
world \emph{is} mathematically: a discrete-time dynamical system whose
state space $\mathcal{S}$ is the product of all bubble state registers,
and whose evolution is governed by a single global update operator.

\begin{definition}[Global World Transition Operator]
  The \emph{world transition operator} is a function
  \[
    \Fworld : \mathcal{S} \to \mathcal{S}
  \]
  mapping the world state at tick $t$ to the world state at tick $t+1$:
  \[
    \world_{t+1} = \Fworld(\world_t).
  \]
\end{definition}

The operator $\Fworld$ is not monolithic. It decomposes into a composition
of domain-specific sub-operators applied in a defined sequence:

\begin{definition}[Compositional Update Order]
  \[
    \Fworld =
    \Fdip \circ \Fcombat \circ \Fecon \circ \Fagent \circ \Feco
  \]
  where:
  \begin{itemize}
    \item $\Feco$ updates ecological fields (forest spread, depletion, climate),
    \item $\Fagent$ updates all agent state machines (movement, production, perception),
    \item $\Fecon$ updates economies (resource accumulation, trade, research),
    \item $\Fcombat$ resolves all combat (damage, capture, unit removal),
    \item $\Fdip$ updates diplomatic state (treaty timers, trust decay, proposals).
  \end{itemize}
\end{definition}

\begin{observation}
  The composition order is not arbitrary. Combat must follow agent movement
  (units cannot fight before they move). Economy must follow agents
  (production depends on current unit positions). Diplomacy follows combat
  (war declarations and peace proposals depend on the tick's battle results).
  The world is a \emph{sequentially composed} dynamical system.
\end{observation}

\begin{theorem}[Non-Commutativity of World Updates]
  In general, $\Fcombat \circ \Fagent \neq \Fagent \circ \Fcombat$.
\end{theorem}
\begin{proof}
  Let agent $A$ be adjacent to agent $B$'s city at tick $t$.
  Under $\Fagent \circ \Fcombat$: combat fires first (no movement has
  occurred), so $A$ attacks from its current position. Under
  $\Fcombat \circ \Fagent$: $A$ moves first, possibly reaching the city
  and attacking from inside its walls. These produce different outcomes.
  $\square$
\end{proof}

This non-commutativity is not a bug: it encodes the \emph{rules of engagement}
of the world. Changing the update order changes which civilization laws hold.
A world where diplomacy is resolved before combat produces different strategic
equilibria than one where combat precedes diplomacy.

\begin{mechanic}{Race Conditions as World Events}
  When two empires simultaneously attempt mutually exclusive actions
  (two settlers claiming the same tile, two armies attacking the same city),
  the world must define a \emph{priority function} over simultaneous events.
  The standard resolution: lower empire ID wins ties. This is an arbitrary
  but deterministic convention---essential for replay integrity.
\end{mechanic}

\section{The Full World Bubble}

\begin{construction}[The Computational World]
  The complete $\world$ bubble contains, in order of construction:
  \begin{enumerate}
    \item \texttt{WorldMap}: $n \times n$ array of \texttt{TileState} bubbles
    \item \texttt{FogSystem}: visibility and exploration tracking
    \item \texttt{EcologyLayer}: forest, depletion, climate
    \item \texttt{EmpireArray}: up to $k$ empire state bubbles, each containing:
      \begin{itemize}
        \item \texttt{CityArray}: up to $m$ city bubbles
        \item \texttt{UnitArray}: up to $p$ unit agent bubbles
        \item \texttt{Treasury}: gold and resource ledger
        \item \texttt{TechTree}: technology state
        \item \texttt{DiplomacyModule}: treaty and trust state
      \end{itemize}
    \item \texttt{EventLog}: append-only history
    \item \texttt{TurnManager}: synchronizes updates across all components
  \end{enumerate}
\end{construction}

\section{The Abstraction Ladder}

The full ascent of the book:
\begin{align*}
  &\text{Signals} \\
  &\quad\to \text{Logic Gates (NAND, AND, OR, NOT, XOR, MUX)} \\
  &\quad\to \text{Arithmetic (half adder, full adder, ALU)} \\
  &\quad\to \text{Memory (bit, register, RAM)} \\
  &\quad\to \text{Sequential State (DFF, latch, counter)} \\
  &\quad\to \text{Spatial Automata (cellular rules, diffusion)} \\
  &\quad\to \text{Tiles (terrain, resources, ownership, visibility)} \\
  &\quad\to \text{Agents (units, cities, empires)} \\
  &\quad\to \text{Economies (food, production, gold, research)} \\
  &\quad\to \text{Expansion (settling, border growth)} \\
  &\quad\to \text{Trade (routes, disruption, logistics)} \\
  &\quad\to \text{Conflict (combat, capture, war)} \\
  &\quad\to \text{Diplomacy (treaties, trust, belief)} \\
  &\quad\to \text{Technology (DAG, prerequisites, unlocks)} \\
  &\quad\to \text{Ecology (depletion, spread, climate)} \\
  &\quad\to \text{History (events, queries, narrative)} \\
  &\quad\to \text{World}
\end{align*}

\section{Victory Conditions as Global Predicates}

A victory condition is a Boolean predicate over the entire world state:

\begin{align*}
  \texttt{domination\_victory} &= \text{(one empire controls all cities)} \\
  \texttt{science\_victory} &= \text{(spaceship launched: technology predicate)} \\
  \texttt{cultural\_victory} &= \text{(cultural influence dominates all tiles)} \\
  \texttt{diplomatic\_victory} &= \text{(elected leader by allied empires)}
\end{align*}

The game ends at the first tick $t^*$ when any victory predicate becomes true.

\section{Worlds That Rewrite Themselves}

The deepest feature of \Sp{} is that bubbles can be created and destroyed
at runtime. A world simulation can spawn new agent bubbles (settlers
founding cities), destroy bubbles (units dying), and modify the rule
structure of existing bubbles (technology unlocking new unit behaviors).

This is not merely parametric variation. It is \emph{structural
self-modification}: the world's computational architecture changes as the
world evolves. History does not just accumulate data. It accumulates
\emph{capabilities}.

\begin{perspective}
  We began with NAND. We end with history.

  The journey from a single gate to a persistent civilization reveals
  something important: complexity does not require a rich initial
  vocabulary. It requires \emph{time} and \emph{iteration}. A world
  is not built from special ingredients. It is built from the repetition
  of simple rules across space and time until the interactions become
  too intricate to trace by hand.

  This is the sense in which a 4X simulation is a philosophical object.
  It shows us, in miniature, how global properties emerge from local
  logic. How sovereignty is nothing but a latch that holds. How history
  is nothing but a trace that accumulates. How civilization is nothing
  but a very large, very patient state machine.

  \Spbox  makes this visible because it refuses to hide the substrate.
  Every empire is a bubble. Every treaty is a shared bubble. Every battle
  is a mutual exclusion race. The world is transparent---and in that
  transparency, educational.
\end{perspective}

\begin{project}{The Complete World}
  \textbf{Final Integration Project.}

  Assemble all components built across Chapters 1--15 into a single
  \Sp{} world simulation. The world must:

  \begin{enumerate}
    \item Initialize a $16 \times 16$ map with procedural terrain.
    \item Spawn 3 empires, each starting with 1 city and 2 units.
    \item Run for 50 turns, updating all subsystems each tick:
      \begin{itemize}
        \item resource accumulation and population growth,
        \item unit movement (rule-based AI),
        \item combat resolution,
        \item diplomatic state,
        \item technology research,
        \item ecological update,
        \item event logging.
      \end{itemize}
    \item Check victory conditions each turn.
    \item Print a full narrative history at game end.
  \end{enumerate}

  \medskip
  \noindent\textit{This project is the capstone of the book.
  A student who completes it has built, from first principles
  and from logic gates upward, a living computational world.}
\end{project}


% ============================================================
% PART VI: PERSISTENCE AND EMERGENCE  (chapters 17-18)
% ============================================================

\part{Persistence and Emergence}

\chapterquote{A world that cannot survive interruption is not a world.
  It is a performance.
}{}

% ── Chapter 17 ───────────────────────────────────────────────
\chapter{Persistence and Save Systems}
\label{ch:persistence}

\chapterquote{Saving is not a technical afterthought.
  It is the world's first act of self-description.
}{}

\begin{background}
  A world that cannot survive interruption is not a world---it is a
  performance. Genuine computational worlds must persist across sessions,
  survive hardware failures, and remain coherent after partial writes.
  These requirements force us to ask a deep question: what \emph{is}
  world state, formally, and how can a finite bitstring faithfully encode
  the entire configuration of a civilization at an arbitrary moment in time?

  Persistence is not merely engineering. It is an ontological commitment.
  When we serialize a world, we claim that the world \emph{is} its state:
  that the full causal history of the simulation is captured in the current
  register contents, event log, and structural layout of bubbles. This
  chapter examines that claim carefully, constructs a save system that
  takes it seriously, and identifies precisely which aspects of world state
  are and are not recoverable from a snapshot.
\end{background}

\section{What Must Be Saved}

The complete state of a 4X world at tick $t$ consists of several layers,
not all of which are equally straightforward to serialize.

\begin{definition}[World Snapshot]
  A \emph{world snapshot} at time $t$ is a tuple
  \[
    \Sigma_t = \bigl(\mathcal{M}_t,\; \mathcal{E}_t,\; \mathcal{A}_t,\;
    \mathcal{L}_t,\; t\bigr)
  \]
  where $\mathcal{M}_t$ is the map state (tile array), $\mathcal{E}_t$ is
  the empire array (all agent states), $\mathcal{A}_t$ is the full event
  log up to tick $t$, $\mathcal{L}_t$ is the set of active latch values
  for all sequential bubbles, and $t$ is the current tick counter.
\end{definition}

The crucial observation is that $\mathcal{M}_t$ and $\mathcal{E}_t$ are
\emph{redundant} with $\mathcal{A}_0$ (the initial conditions) and the
complete event log $\mathcal{A}_t$ if the simulation is deterministic.
In principle, one could save only the initial seed and the event log, then
replay to reconstruct any later state. In practice, this is prohibitively
slow for long games; thus real save systems maintain both the current
state \emph{and} the historical log, trading storage for recovery speed.

\section{Serialization as World Projection}

\begin{definition}[Serialization]
  A \emph{serialization function} is a map
  \[
    \mathrm{ser} : \Sigma \to \{0,1\}^*
  \]
  assigning a finite binary string to each world snapshot, together with
  a left-inverse
  \[
    \mathrm{deser} : \{0,1\}^* \to \Sigma \cup \{\bot\}
  \]
  satisfying $\mathrm{deser}(\mathrm{ser}(\sigma)) = \sigma$ for all
  $\sigma \in \Sigma$.
\end{definition}

\begin{theorem}[Serialization Existence]
  For any finite world $\world$ with bounded state registers, a
  serialization function exists and can be constructed by recursive
  structural traversal of the bubble tree.
\end{theorem}
\begin{proof}
  By structural induction on the bubble hierarchy. The base case is a
  primitive state register of width $w$: its serialization is the $w$-bit
  register contents, and deserialization is the identity. For a composite
  bubble $B$ with children $B_1, \ldots, B_k$, define
  \[
    \mathrm{ser}(B) = \mathrm{ser}(B_1) \,\|\, \mathrm{ser}(B_2)
    \,\|\, \cdots \,\|\, \mathrm{ser}(B_k)
  \]
  where $\|$ denotes concatenation. Since the structure of $B$ is fixed
  (compiled into the simulation code, not stored in state),
  $\mathrm{deser}$ knows exactly how many bits to allocate to each child,
  so it can reconstruct each $B_i$ independently and assemble them.
  By induction, each $\mathrm{deser}(B_i)$ succeeds, so
  $\mathrm{deser}(B)$ succeeds. The world $\world$ is the root bubble;
  apply the composite case.
\end{proof}

\section{Rollback and Replay}

\begin{defbox}{Rollback}
  A \emph{rollback} to tick $t' < t$ is the operation of deserializing
  the snapshot $\Sigma_{t'}$, discarding all events with timestamp
  $> t'$, and resuming execution from $\Sigma_{t'}$.
\end{defbox}

\begin{mechanic}{Replay Integrity}
  A simulation supports \emph{replay integrity} if, for any initial
  snapshot $\Sigma_0$ and event sequence $\mathcal{A}_t$,
  re-executing $\mathcal{A}_t$ from $\Sigma_0$ produces the same
  $\Sigma_t$ as the original run. This requires:
  (1) deterministic bubble evaluation (no hidden randomness outside
  the event log);
  (2) complete event capture (every state-modifying action is logged);
  (3) monotonic time (events are applied in order).
\end{mechanic}

\begin{implementation}
\begin{lstlisting}[caption={Snapshot and restore in \Spbox }]
bubble SaveSystem {
    inputs:  world_state, command[2]  -- 00=idle 01=save 10=load 11=replay

    latch snapshot[MAX_STATE_BITS]
    latch event_log[MAX_EVENTS][64]
    latch log_head[16]

    when command = SAVE do
        snapshot    := serialize(world_state)
    end

    when command = LOAD do
        world_state := deserialize(snapshot)
        log_head    := snapshot.tick
    end

    when command = REPLAY do
        world_state := deserialize(snapshot)
        let i = 0
        while i < log_head do
            world_state := apply(event_log[i], world_state)
            i = i + 1
        end
    end
}
\end{lstlisting}
\end{implementation}

\begin{perspective}
  The ability to serialize and replay a world reveals something
  philosophically remarkable: the world is a \emph{function}, not a
  substance. It has no hidden essences, no unobserved interior that
  escapes the snapshot. Everything that the world \emph{is} at time $t$
  is captured in $\Sigma_t$. This is the computational correlate of
  Hume's bundle theory of identity: the world is nothing over and above
  the bundle of its current states and their history of production.

  Contrast this with the Aristotelian view, in which substances have
  natures that transcend their current properties~\cite{Aristotle}.
  The serializable world has no such nature: it is entirely transparent
  to its own description.
\end{perspective}

\begin{project}{Complete Save/Load Architecture}
  \begin{enumerate}
    \item Implement \texttt{serialize} and \texttt{deserialize} for the
          full \texttt{WorldBubble} from Chapter~16.
    \item Add a \texttt{CheckpointManager} that saves automatically
          every 10 turns.
    \item Implement rollback to any prior checkpoint.
    \item Build a replay viewer: given a saved game, replay and print
          the event log as a chronicle.
    \item Verify replay integrity: confirm that replaying from the
          initial state produces the same final state as the original run.
  \end{enumerate}
\end{project}

% ── Chapter 18 ───────────────────────────────────────────────
\chapter{Emergence}
\label{ch:emergence}

\chapterquote{More is different.
}{--- Philip W.\ Anderson}

\begin{background}
  We have now constructed every layer of the 4X world individually.
  In this chapter, we step back and study the world as a whole. When
  the layers interact---when ecology constrains resources, resources
  constrain population, population drives military production, military
  production enables expansion, and expansion modifies ecology---new
  behaviors appear that were not written into any individual bubble.
  This is emergence: the arising of global patterns from local rules.

  Emergence is not mystical. It is a precise phenomenon with a
  mathematical definition, and it has a specific relationship to the
  architecture of \Spbox  that we can state formally. Understanding
  emergence is the final step before we turn to the philosophical
  foundations of the Spherepop framework itself.
\end{background}

\section{Formal Definition of Emergence}

\begin{definition}[Emergent Property]
  A property $P$ of a world $\world$ is \emph{emergent} with respect to
  a partition $\{L_1, \ldots, L_k\}$ of $\world$'s bubbles into layers if:
  (1) $P(\world)$ is true, and
  (2) $P$ is not expressible as a Boolean combination of properties
  of any single layer $L_i$ in isolation.
\end{definition}

The canonical 4X emergent properties are empire formation, economic
collapse, migration pressure, cultural divergence, power-law city sizes,
and phase transitions---each arising from the simultaneous interaction
of multiple layers, none of which produces the phenomenon on its own.

\section{Power Laws and Zipf's Law}

In historical civilizations, city populations obey a power law: if cities
are ranked by population, the $k$-th city has population approximately
$P_1 / k$, where $P_1$ is the largest city's population. This is Zipf's Law.

\begin{theorem}[Zipf Emergence Condition]
  In a 4X world with preferential attachment in trade route formation,
  the steady-state city population distribution follows a power law
  with exponent $\alpha \approx 1$.
\end{theorem}
\begin{proof}[Proof sketch]
  Let $C_i(t)$ be the population of city $i$ at tick $t$, and let the
  growth rate be proportional to trade route degree $d_i(t)$:
  $C_i(t+1) = C_i(t) + \lambda d_i(t)$.
  Under preferential attachment, a new trade route connects to city $i$
  with probability $d_i / \sum_j d_j$. This is the Barab\'{a}si--Albert
  mechanism, whose steady-state degree distribution is
  $P(d) \propto d^{-3}$. Since $C \propto d$, the rank distribution
  satisfies Zipf's Law with exponent 1. $\square$
\end{proof}

\section{Phase Transitions in Civilization}

\begin{definition}[Civilizational Phase Transition]
  A \emph{phase transition} in world $\world$ occurs at tick $t^*$ when
  a global order parameter $\Phi(t)$ undergoes a discontinuous jump:
  \[
    \lim_{t \to t^{*-}} \Phi(t) \neq \lim_{t \to t^{*+}} \Phi(t).
  \]
\end{definition}

\begin{mechanic}{Percolation and Empire Collapse}
  Model imperial territory as a graph $G = (V, E)$ where $V$ is the set
  of owned tiles and $E$ connects adjacent owned tiles. When military
  defeats remove tiles, connectivity can drop catastrophically at a
  threshold removal fraction $p_c \approx 0.407$ (the site percolation
  threshold for a square lattice). Below $p_c$, the empire fragments into
  independent enclaves---a phase transition from unified to fractured.
\end{mechanic}

\begin{project}{Emergence Analysis}
  \begin{enumerate}
    \item Run 5 complete 100-turn simulations with 4 empires each.
    \item Plot city population distributions; test fit to power law.
    \item Measure territorial connectivity of each empire at each turn.
    \item Identify turns where connectivity drops sharply.
    \item Record cultural divergence: measure Hamming distance between
          technology vectors of different empires over time.
  \end{enumerate}
\end{project}

% ============================================================
% PART VII: HISTORY BEFORE STATE
% ============================================================

\part{History Before State}

\chapterquote{Objects are not primary. Histories are primary.
}{}

% ── Chapter 19 ───────────────────────────────────────────────
\chapter{Events Before Objects}
\label{ch:events}

\chapterquote{The present moment always will have been.
}{--- after Wittgenstein~\cite{Wittgenstein}}

\begin{background}
  Most computational systems begin with objects: data structures
  occupying memory, accessed by name, modified in place. A city is a
  record. A tile is an array cell. An empire is a struct. This is the
  \emph{state-centric} view of computation, and it is so natural to
  programmers that its presuppositions are rarely examined.

  Spherepop begins elsewhere. It takes seriously the observation,
  developed in process philosophy by Whitehead~\cite{Whitehead} and in
  the philosophy of events by Kim~\cite{Kim}, that \emph{events are
  ontologically prior to objects}. A city is not a data structure that
  occasionally changes. A city is the accumulated history of all events
  that brought it into being and sustained it: the founding event, the
  population growth events, the construction events, the siege events.
  Remove the history and you do not have a city at rest---you have
  nothing at all. Identity is causal, not structural.

  This chapter constructs the event-sourcing substrate that underlies the
  Spherepop philosophy, proves its equivalence to the state-based model
  developed in Chapters~3--5, and then demonstrates why the event model
  is strictly \emph{richer}: it preserves information that the state model
  discards.
\end{background}

\section{The Event-Centric World Model}

\begin{definition}[Event]
  An \emph{event} is an irreversible world operation. Formally, an event
  $\varepsilon$ is a morphism
  \[
    \varepsilon : X \to X'
  \]
  where $X$ and $X'$ are world option spaces (sets of admissible future
  states), with the constraint that $X' \subseteq X$ (events never expand
  the option space).
\end{definition}

\begin{definition}[World History]
  The \emph{history} of a world from initial configuration $X_0$ is the
  composite morphism
  \[
    H = \varepsilon_n \circ \varepsilon_{n-1} \circ \cdots \circ \varepsilon_1
    : X_0 \to X_n.
  \]
\end{definition}

\section{Event Sourcing as Information Preservation}

\begin{theorem}[Events Preserve More Than State]
  Distinct histories can produce identical present states, but an event
  log distinguishes them while a state snapshot does not.
\end{theorem}
\begin{proof}
  Consider two histories on a minimal world with one tile $T$ having
  ownership register $O \in \{0,1\}$:
  \begin{align*}
    \mathcal{L}_1 &= (\texttt{Claim}(T, e_1)) \\
    \mathcal{L}_2 &= (\texttt{Claim}(T, e_1),\;
                      \texttt{Release}(T),\; \texttt{Claim}(T, e_1)).
  \end{align*}
  Both produce final state $O = 1$. A snapshot contains only $O = 1$;
  it cannot distinguish $\mathcal{L}_1$ from $\mathcal{L}_2$. But the
  logs differ: $|\mathcal{L}_1| = 1 \neq 3 = |\mathcal{L}_2|$. The
  diplomatic consequence (empire $e_1$ has a history of releasing
  territory) is recoverable from $\mathcal{L}_2$ but from neither
  snapshot. $\square$
\end{proof}

\section{Historical Identity}

\begin{definition}[Historical Identity]
  Two entities $A$ and $B$ are \emph{historically identical} if and only
  if their event logs are identical:
  \[
    A \equiv_H B \;\Longleftrightarrow\; \mathcal{L}(A) = \mathcal{L}(B).
  \]
\end{definition}

\begin{corollary}
  Historical identity is strictly finer than state identity:
  $A \equiv_H B \Rightarrow \mathrm{state}(A) = \mathrm{state}(B)$,
  but the converse fails.
\end{corollary}

\begin{implementation}
\begin{lstlisting}[caption={Append-only event sourcing in \Spbox }]
bubble EventSourcedWorld {
    latch event_log[MAX_EVENTS][128]
    latch log_head[16]

    bubble QueryOwner {
        inputs:  x[8], y[8]
        output:  owner[4]

        let i = log_head
        owner = 0
        while i > 0 do
            i = i - 1
            if event_log[i].type = CLAIM
               AND event_log[i].tile_x = x
               AND event_log[i].tile_y = y then
                owner = event_log[i].actor
                i = 0
            end
            if event_log[i].type = RELEASE
               AND event_log[i].tile_x = x
               AND event_log[i].tile_y = y then
                owner = 0
                i = 0
            end
        end
    }

    bubble AppendEvent {
        inputs:  event[128]
        when log_head < MAX_EVENTS do
            event_log[log_head] := event
            log_head := log_head + 1
        end
    }
}
\end{lstlisting}
\end{implementation}

\begin{perspective}
  The event-sourcing architecture embodies Whitehead's process
  philosophy~\cite{Whitehead}: the fundamental furniture of the universe
  consists not of enduring substances but of momentary occasions of
  experience, each of which inherits and transforms the actual world of
  its causal ancestors. The append-only log is the formal analogue of
  Whitehead's creative advance: once an event is written, it cannot be
  unwritten. The past is immutable; only the future remains open.
\end{perspective}

\begin{project}{Event-Sourced Tile Ownership}
  \begin{enumerate}
    \item Replace all mutable \texttt{owner} registers with an event log.
    \item Implement \texttt{QueryOwner} by log scan.
    \item Implement \texttt{QueryHistory}: return all ownership events
          for a given tile.
    \item Implement \texttt{TrustFromHistory}: count violated treaties
          in the log to compute diplomatic trust.
    \item Design a hybrid cache: register snapshot for speed plus full
          log for auditability.
  \end{enumerate}
\end{project}

% ── Chapter 20 ───────────────────────────────────────────────
\chapter{Nested Worlds and Bubble Scope}
\label{ch:nesting}

\chapterquote{A battle exists inside a war.
  A war exists inside a civilization.
  A civilization exists inside a planetary history.
}{}

\begin{background}
  Computation is fundamentally about scope: what a computation can see,
  what it can modify, and what is hidden from it. In lambda
  calculus~\cite{Barendregt}, scope is defined by $\lambda$-abstraction.
  In \Spbox , scope is defined by \emph{containment}: a bubble's interior
  is hidden from everything outside its boundary. A bubble boundary is
  a \emph{distinction}---it separates the world into inside and outside.

  For world simulation, nested scope has immediate strategic significance.
  A city's internal economy is hidden from the empire's strategic AI,
  which sees only city outputs. The empire's military strategy is hidden
  from the world's ecological processes, which see only physical effects.
  Nested scope is what allows a world to be simultaneously local and
  global: each bubble is complete in itself and also a component of
  something larger.
\end{background}

\section{The Scope Hierarchy}

\begin{definition}[Scope Hierarchy]
  A \emph{scope hierarchy} is a rooted tree $\mathcal{T} = (N, E, r)$
  where each node $n \in N$ is a bubble, edges represent containment,
  and $r$ is the root (the world).
\end{definition}

\begin{definition}[Lexical Scope in Spherepop]
  A bubble $B$ has \emph{lexical access} to its own internal registers,
  the input ports declared by its parent, and any bubble explicitly
  passed as an argument. $B$ does not have access to the internals of
  its siblings or ancestors beyond declared ports.
\end{definition}

\begin{theorem}[Scope Compositionality]
  In a scope hierarchy $\mathcal{T}$, if each bubble correctly implements
  its specification independently of all sibling and ancestor internals,
  then any composite bubble formed by connecting bubbles in $\mathcal{T}$
  also correctly implements its specification.
\end{theorem}
\begin{proof}
  By structural induction on $\mathcal{T}$. Base case: a leaf
  (primitive) bubble is correct by assumption. For an interior node $B$
  with children $B_1, \ldots, B_k$: by the induction hypothesis, each
  $B_i$ correctly implements its specification. The wiring of $B$
  connects output ports of some children to input ports of others,
  respecting interface contracts. Since each $B_i$'s behavior depends
  only on its declared inputs (by lexical scope) and those inputs are
  correctly produced (by induction), $B$'s output is correctly
  determined. $\square$
\end{proof}

\begin{implementation}
\begin{lstlisting}[caption={Nested city-district economy}]
bubble World {
    bubble Empire {
        bubble City {
            bubble District {
                latch commerce[8]
                bubble Market {
                    inputs:  supply[8], demand[8]
                    output:  price[8]
                    price = demand / (supply + 1)
                }
            }
            -- City sees district outputs only
            inputs:   district_output[16]
            output:   city_gold[8], city_food[8]
        }
    }
}
\end{lstlisting}
\end{implementation}

\begin{perspective}
  Nested scope is the computational realization of Leibniz's monads:
  each bubble is a ``windowless'' computation that sees the world only
  through its declared ports. This enforced opacity is the source of
  modularity. Without scope, every bubble would need to know everything,
  and the simulation would collapse into a single undifferentiated
  process. With scope, each layer can be designed, verified, and replaced
  independently---and the whole can be greater than the sum of its parts.
\end{perspective}

\begin{project}{Nested Local Economies}
  \begin{enumerate}
    \item Implement a three-level hierarchy:
          \texttt{World} $\supset$ \texttt{Empire} $\supset$
          \texttt{City} $\supset$ \texttt{District}.
    \item Each district manages labor allocation among farming, mining,
          and commerce.
    \item Cities see only district output totals.
    \item Empires see only city output totals.
    \item Implement cross-city trade at the empire level accessing only
          declared city output ports.
  \end{enumerate}
\end{project}

% ── Chapter 21 ───────────────────────────────────────────────
\chapter{The Pop Operator}
\label{ch:pop}

\chapterquote{To choose is to exclude. The value of a decision lies precisely
  in what it forecloses.
}{}

\begin{background}
  The name ``Spherepop'' is not decorative. The \emph{pop} is the
  primitive act of the system: the irreversible commitment that collapses
  an ambiguous set of futures into a single actualized one. Before a
  Settler unit settles, a city could arise anywhere. After the settlement
  event fires, one location is actualized and all others are permanently
  excluded. The pop is the moment of decision.

  Kierkegaard's ``either/or,'' Sartre's radical freedom, Heidegger's
  \emph{Entschlossenheit}, and Badiou's event~\cite{Badiou} all circle
  the same structure: the moment at which the indeterminate becomes
  determinate, irreversibly. What Spherepop contributes is a precise
  formal language for this moment, grounded in the mathematics of option
  spaces and contraction operators.

  The pop is also the moment where games acquire meaning. A game in which
  no decision is irreversible is a game without stakes. The sense of
  weight that a well-designed 4X game produces---the feeling that founding
  this city here, in this moment, with these resources, commits you to a
  trajectory from which there is no return---is the phenomenological
  signature of the pop operator.
\end{background}

\section{Option Spaces and Contraction}

\begin{definition}[Option Space]
  An \emph{option space} $\Omega$ is a set of admissible future world
  trajectories. Each element $\omega \in \Omega$ is a maximal consistent
  extension of the current world state to all future ticks.
\end{definition}

\begin{definition}[Pop Operator]
  A \emph{pop} is a function $\mathbf{P}_\varepsilon : \Omega \to \Omega'$
  associated to an event $\varepsilon$, satisfying:
  (1)~\textbf{Contraction:} $\Omega' \subseteq \Omega$.
  (2)~\textbf{Irreversibility:} There is no operator $\mathbf{Q}$ such
  that $\mathbf{Q} \circ \mathbf{P}_\varepsilon = \mathrm{id}_\Omega$.
  (3)~\textbf{Commitment:} Every $\omega' \in \Omega'$ is consistent with
  $\varepsilon$ having occurred.
\end{definition}

\begin{theorem}[Strict Contraction]
  If event $\varepsilon$ is non-trivial, then $|\Omega'| < |\Omega|$.
\end{theorem}
\begin{proof}
  Let $\omega^* \in \Omega$ be a trajectory inconsistent with $\varepsilon$
  (exists by non-triviality). Then $\omega^* \notin \Omega'$ by the
  commitment condition, so $\Omega' \subseteq \Omega \setminus \{\omega^*\}$,
  giving $|\Omega'| \leq |\Omega| - 1 < |\Omega|$. $\square$
\end{proof}

\begin{corollary}[Historical Entropy is Non-Increasing]
  Define $S(\Omega) = \log_2 |\Omega|$.
  For any non-trivial pop, $S(\Omega') < S(\Omega)$.
\end{corollary}

This connects the Spherepop pop to Landauer's principle~\cite{Landauer}:
every irreversible bit-erasure has a thermodynamic cost.

\begin{implementation}
\begin{lstlisting}[caption={Pop operator for city settlement}]
bubble PopSettle {
    inputs:
        empire_id[4], tile_x[8], tile_y[8], settler_id[8]

    let allowed = SettlePermission(empire_id, tile_x, tile_y)

    when allowed = 1 do
        -- 1. Remove settler (irreversible)
        world.units[settler_id].alive := 0
        -- 2. Found city (irreversible event)
        AppendEvent(FOUND_CITY, empire_id, tile_x, tile_y, tick)
        -- 3. Claim tile (irreversible)
        AppendEvent(CLAIM, empire_id, tile_x, tile_y, tick)
    end
}
\end{lstlisting}
\end{implementation}

\begin{perspective}
  The pop operator formalizes what game designers call ``meaningful
  choice.'' A choice is meaningful precisely when it irreversibly
  constrains the future. Games without irreversible decisions---where
  every action can be undone by the save/load cycle---have lower
  strategic depth not because they lack complexity but because they lack
  \emph{commitment}. The pop is the computational measure of decisional
  weight. Badiou~\cite{Badiou} argues that an event is precisely that
  which ruptures the existing order and forces a new fidelity. The
  settlement pop is a miniature Badiousian event.
\end{perspective}

\begin{project}{Irreversible Settlement System}
  \begin{enumerate}
    \item Implement the full \texttt{PopSettle} bubble with precondition
          checking and event logging.
    \item Build a \texttt{PopDeclareWar} bubble that excludes all
          trajectories of peaceful coexistence with the target.
    \item Build a \texttt{PopChooseTech}: committing research to one
          technology temporarily forecloses alternatives.
    \item Measure the option space reduction (entropy decrease) for each
          pop type by counting reachable states before and after.
  \end{enumerate}
\end{project}

% ── Chapter 22 ───────────────────────────────────────────────
\chapter{Merge and Collapse}
\label{ch:merge}

\begin{background}
  If the pop contracts an option space by committing to a specific event,
  \emph{merge} operates on entities within the world: combining two
  structures into one, collapsing distinctions that were previously
  maintained. Empires merge territories. Cultures merge vocabularies.
  Trade networks merge into single markets. These are not additions;
  they are \emph{compressions}---operations that identify previously
  distinct entities and represent their union as a single canonical form.

  The literature on Conflict-free Replicated Data Types (CRDTs)
  in computer science~\cite{Meijer} formalizes this idea: a data
  structure supports merge if it forms a join-semilattice. Spherepop's
  merge operator inherits this structure and extends it to the richer
  domain of world entities.
\end{background}

\section{The Merge Semilattice}

\begin{definition}[Merge Semilattice]
  A set $\mathcal{R}$ with binary operator $\mu$ is a
  \emph{merge semilattice} if $\mu$ is commutative, associative,
  and idempotent.
\end{definition}

\begin{theorem}[Territory Merge is a Semilattice]
  Let $\mathcal{R} = \mathcal{P}(\text{Tiles})$ with $\mu(A,B) = A \cup B$.
  Then $(\mathcal{R}, \mu)$ is a merge semilattice.
\end{theorem}
\begin{proof}
  Set union is commutative ($A \cup B = B \cup A$), associative
  ($(A \cup B) \cup C = A \cup (B \cup C)$), and idempotent
  ($A \cup A = A$). $\square$
\end{proof}

\section{Collapse and Canonicalization}

\begin{definition}[Collapse Operator]
  A \emph{collapse operator} $\kappa : \mathcal{R} \to \widehat{\mathcal{R}}$
  maps each entity to its canonical representative, satisfying:
  \[
    \kappa(\mu(A,B)) = \kappa(\mu(\kappa(A), \kappa(B))).
  \]
\end{definition}

\begin{example}[Cultural Assimilation]
  Empire $e_1$ has culture set $C_1 = \{\text{Latin}, \text{Greek}\}$
  and empire $e_2$ has $C_2 = \{\text{Greek}, \text{Aramaic}\}$.
  Cultural merge: $C_{\mathrm{merged}} = \{\text{Latin}, \text{Greek},
  \text{Aramaic}\}$. Collapse: $\kappa$ selects the dominant culture
  (highest population weight), producing a canonical representative.
\end{example}

\begin{project}{Cultural Assimilation Mechanics}
  \begin{enumerate}
    \item Define culture as a Boolean vector over 16 cultural traits.
    \item Implement \texttt{CultureMerge} as set union on trait vectors.
    \item Implement \texttt{CultureCollapse}: return dominant culture.
    \item Build \texttt{AssimilationPressure}: tiles within a cultural
          sphere gradually adopt dominant traits.
    \item Simulate two civilizations meeting and assimilating over 30 turns.
  \end{enumerate}
\end{project}

% ── Chapter 23 ───────────────────────────────────────────────
\chapter{History as Identity}
\label{ch:histident}

\begin{background}
  In Chapter~19, we established that events are ontologically prior to
  objects. In Chapter~21, we showed that events irreversibly constrain
  the future. Now we draw out the identity-theoretic consequences: if
  what an entity \emph{is} is its history, then two entities are the same
  only if they share the same history. This is a radical thesis---far
  stronger than extensional identity (same current properties) or modal
  identity (same counterfactual profile). But it is the thesis that
  Spherepop's event-sourcing architecture implicitly commits to, and it
  has significant consequences for persistence, replication, and memory
  in simulations.
\end{background}

\section{Provenance and Lineage}

\begin{definition}[Provenance]
  The \emph{provenance} of entity $A$ at time $t$ is the ordered sequence
  of all events that produced $A$:
  \[
    \mathrm{prov}(A, t) = (\varepsilon_1^A, \varepsilon_2^A, \ldots,
    \varepsilon_{n(t)}^A).
  \]
\end{definition}

\begin{theorem}[Provenance Determines State]
  In a deterministic simulation, $\mathrm{prov}(A,t)$ together with the
  initial conditions $X_0$ uniquely determines $\mathrm{state}(A,t)$.
\end{theorem}
\begin{proof}
  By induction on $n(t)$. Base case $n(t) = 0$: state equals $A$'s
  initial state, determined by $X_0$. Inductive step: by hypothesis
  $\mathrm{prov}(A,t')$ determines $\mathrm{state}(A,t')$ for the tick
  $t'$ preceding event $\varepsilon_{k+1}^A$, and the deterministic
  transition function $\delta$ maps
  $(\mathrm{state}(A,t'), \varepsilon_{k+1}^A)$ to $\mathrm{state}(A,t)$.
  $\square$
\end{proof}

\begin{implementation}
\begin{lstlisting}[caption={Timeline reconstruction}]
bubble TimelineReconstructor {
    inputs:  entity_id[8], target_tick[16]
    output:  reconstructed_state[512]

    let state = initial_state(entity_id)
    let i = 0
    while i < global_log_head do
        if event_log[i].actor = entity_id
           OR event_log[i].target = entity_id then
            if event_log[i].tick <= target_tick then
                state = apply_event(state, event_log[i])
            end
        end
        i = i + 1
    end
    reconstructed_state = state
}
\end{lstlisting}
\end{implementation}

\begin{project}{Historical Replay and Timeline Reconstruction}
  \begin{enumerate}
    \item Implement \texttt{TimelineReconstructor} for cities.
    \item Build a \texttt{CityBiography}: given a city ID, produce a
          human-readable chronicle of all events in the city's history.
    \item Implement lineage tracking: captured and refounded cities
          inherit the old city's event log as a prefix.
    \item Build a divergence metric: compute the tick at which two
          empires first diverged and the Hamming distance between
          their current technology and culture vectors.
  \end{enumerate}
\end{project}

% ============================================================
% PART VIII: GEOMETRIC COMPUTATION
% ============================================================

\part{Geometric Computation}

\chapterquote{Space is not a container. It is a constraint structure.
}{}

% ── Chapter 24 ───────────────────────────────────────────────
\chapter{Geometry as Computation}
\label{ch:geometry}

\begin{background}
  Throughout this book, geometry has appeared implicitly: in tile
  adjacency, vision radii, path-checking of trade routes, and
  percolation thresholds of territorial connectivity. We now make
  geometry explicit. The spatial structure of the world is not merely
  a rendering concern---it is a computational structure that determines
  what operations are possible, what their costs are, and what emergent
  patterns can arise.

  In the Spherepop framework, geometry is understood as a
  \emph{constraint structure on option spaces}. Mountains are not merely
  visually distinct from plains---they constrain movement, visibility,
  and resource extraction in ways that structurally shape the space of
  possible histories.
\end{background}

\section{Tile Graphs and Spatial Topology}

\begin{definition}[Tile Graph]
  The \emph{tile graph} of a world is the undirected graph $G = (V, E)$
  where $V$ is the set of passable tiles and $(u, v) \in E$ iff $u$ and
  $v$ are adjacent.
\end{definition}

\begin{theorem}[Connectivity and Reachability]
  A unit at position $u$ can reach position $v$ in at most $k$ turns
  (with movement cost 1 per tile) if and only if the graph distance
  $d_G(u, v) \leq k$.
\end{theorem}
\begin{proof}
  The unit's reachable set after $k$ turns is exactly the ball of
  radius $k$ in the tile graph: $B_k(u) = \{v : d_G(u,v) \leq k\}$.
  This follows from BFS on $G$. $\square$
\end{proof}

\section{Region-Based Territorial Reasoning}

\begin{defbox}{Territorial Compactness}
  The \emph{compactness} of a territory $T$ is
  \[
    \mathrm{compact}(T) = \frac{|V(T)|^2}{|E(T)|}.
  \]
  High compactness means short internal travel distances; low compactness
  means a fragile, extended territory.
\end{defbox}

\begin{project}{Region-Based Territorial Reasoning}
  \begin{enumerate}
    \item Implement \texttt{ConnectedComponents} for each empire.
    \item Implement \texttt{Compactness} per turn.
    \item Build \texttt{CutVertexDetector}: tiles whose removal
          disconnects an empire's territory.
    \item Visualize compactness and connectivity over a 50-turn simulation.
  \end{enumerate}
\end{project}

% ── Chapter 25 ───────────────────────────────────────────────
\chapter{Optionality and Entropy}
\label{ch:optionality}

\begin{background}
  Chapter~21 introduced the pop as a contraction of the option space.
  Now we develop optionality as a \emph{positive} quantity---a measure
  of a civilization's strategic power. A civilization with high
  optionality has many accessible futures; one with low optionality
  is locked in. This connects Spherepop to the theory of real options
  in economics and to Shannon's information theory~\cite{Shannon}: the
  same formula measures uncertainty in all contexts.
\end{background}

\section{The Optionality Field}

\begin{definition}[Optionality]
  The \emph{optionality} of civilization $C$ at tick $t$ is
  \[
    \Phi_C(t) = \log_2 |\mathcal{T}_C(t)|
  \]
  where $\mathcal{T}_C(t)$ is the set of all trajectories consistent
  with $C$'s event history up to $t$.
\end{definition}

In practice we decompose this into domain-specific components:

\begin{definition}[Accessibility Field]
  \[
    \Phi_C(t) = \sum_{d} w_d \log_2 A_d(t)
  \]
  where $A_d(t)$ counts accessible actions in domain $d$ (military,
  economic, scientific, ecological, diplomatic) and $w_d$ are weights.
\end{definition}

\begin{theorem}[Civilizational Collapse Criterion]
  A civilizational collapse occurs at tick $t^*$ if
  $\frac{d}{dt}\Phi_C\big|_{t=t^*} < -\theta$ for some threshold
  $\theta > 0$, sustained over a window $\Delta t$: rapid, sustained
  loss of accessible futures.
\end{theorem}
\begin{proof}[Derivation]
  The discrete derivative $\Delta\Phi_C(t) = \Phi_C(t+1) - \Phi_C(t)$
  captures per-turn change. Collapse is the regime where adverse
  events dominate: the civilization loses more future options per turn
  than it gains. The threshold $\theta$ calibrates severity; $\Delta t$
  prevents classifying temporary setbacks as collapses. $\square$
\end{proof}

\begin{project}{Optionality Metrics for AI Decision-Making}
  \begin{enumerate}
    \item Implement \texttt{OptionalityMeter} for all four empires.
    \item Record optionality at each turn over a 50-turn simulation.
    \item Build a \texttt{MaxOptionalityAI}: select actions that
          maximize $\Phi$ at the next tick.
    \item Compare against a greedy resource-maximizing AI over 20 runs.
    \item Identify how far in advance $\Phi$ begins declining before
          elimination.
  \end{enumerate}
\end{project}

% ── Chapter 26 ───────────────────────────────────────────────
\chapter{Trajectory Collapse}
\label{ch:trajcollapse}

\begin{background}
  The previous chapter measured optionality as a scalar. This chapter
  develops the geometric structure of trajectory space: how ambiguous
  futures form a probability simplex, how strategic uncertainty is
  distributed over it, and how the pop operator effects a geometric
  collapse. This connects Spherepop to Bayesian inference and Shannon
  entropy~\cite{Shannon}: the same formula---
  $H = -\sum p_i \log p_i$---measures uncertainty in all three contexts.
\end{background}

\section{Hypothesis Manifolds}

\begin{definition}[Trajectory Entropy]
  Given scenarios $\mathcal{H} = \{h_1, \ldots, h_n\}$ with probabilities
  $p_i$, the \emph{trajectory entropy} is:
  \[
    H_t = -\sum_{i=1}^n p_i \log_2 p_i.
  \]
\end{definition}

\begin{theorem}[Pop Reduces Trajectory Entropy]
  For any non-trivial pop eliminating scenario $h_k$:
  \[
    H_{t+1} \leq H_t
  \]
  with equality iff $p_k = 0$.
\end{theorem}
\begin{proof}
  Let $p' = (p_1', \ldots, p_n')$ be the post-pop distribution:
  $p_k' = 0$ and $p_i' = p_i/(1-p_k)$ for $i \neq k$.
  By the data-processing inequality, eliminating an outcome and
  renormalizing never increases entropy: $H' \leq H$. Equality holds
  iff $p_k = 0$ (the pop changes nothing). $\square$
\end{proof}

\begin{project}{Branching Historical Simulations}
  \begin{enumerate}
    \item Implement a \texttt{BranchingSimulator}: at decision points,
          run $k$ parallel branches with different choices.
    \item Measure trajectory entropy at each branch point.
    \item Build a \texttt{ScenariosDatabase}: store and compare final
          states of all branches.
    \item Visualize the branching tree annotated with entropy at each node.
    \item Identify which decisions cause the largest entropy reduction.
  \end{enumerate}
\end{project}

% ============================================================
% PART IX: CATEGORY THEORY AND WORLD CONSTRUCTION
% ============================================================

\part{Category Theory and World Construction}

\chapterquote{A morphism is not a transformation of objects.
  It is a relationship that composes.
}{--- after Mac Lane~\cite{MacLane}}

% ── Chapter 27 ───────────────────────────────────────────────
\chapter{Morphisms and Historical Worlds}
\label{ch:morphisms}

\begin{background}
  Category theory~\cite{MacLane} is the mathematics of structured
  composition. It asks not ``what is this thing?'' but ``how does this
  thing relate to other things of the same kind?'' For world simulation,
  we want to understand transformations between worlds---how a world at
  tick $t$ relates to the world at tick $t+1$, how one empire's history
  relates to another's, how a local economy relates to the global one.

  In this chapter, we construct the category $\mathbf{Hist}$ whose
  objects are option spaces and whose morphisms are events. We show that
  historical composition is categorical composition, derive consequences
  for victory conditions and technology trees, and explore the connection
  to Girard's linear logic~\cite{Girard}---a logic where resources are
  consumed rather than duplicated, mirroring the irreversibility of
  in-game actions.
\end{background}

\section{The History Category}

\begin{definition}[The Category $\mathbf{Hist}$]
  $\mathbf{Hist}$ has option spaces as objects; morphisms
  $\mathrm{Hom}(\Omega_i, \Omega_j)$ are events
  $\varepsilon : \Omega_i \to \Omega_j$; composition is sequential
  occurrence; identity is the null event.
\end{definition}

\begin{theorem}[$\mathbf{Hist}$ is a Category]
  The structure $\mathbf{Hist}$ satisfies the category axioms.
\end{theorem}
\begin{proof}
  \textit{Identity:} $\varepsilon \circ \mathrm{id}_{\Omega_0} = \varepsilon$
  and $\mathrm{id}_{\Omega_1} \circ \varepsilon = \varepsilon$ because
  the null event leaves the option space unchanged.
  \textit{Associativity:} Both $(\varepsilon_3 \circ \varepsilon_2)
  \circ \varepsilon_1$ and $\varepsilon_3 \circ (\varepsilon_2 \circ
  \varepsilon_1)$ apply events in order $\varepsilon_1, \varepsilon_2,
  \varepsilon_3$; sequential composition is associative. $\square$
\end{proof}

\begin{theorem}[Victory Condition as Natural Transformation]
  A monotone victory condition $V$ is a natural transformation
  $V : F \Rightarrow \mathbf{2}$ from the state functor to the constant
  functor on $\{0,1\}$.
\end{theorem}
\begin{proof}[Proof sketch]
  At each option space $\Omega$, $V_\Omega : F(\Omega) \to \{0,1\}$
  evaluates the victory predicate. Naturality (once true, stays true)
  requires $V_{\Omega'} \circ F(\varepsilon) = V_\Omega$ for all
  events $\varepsilon$. This holds when victory is terminal---precisely
  the design criterion for valid victory conditions. $\square$
\end{proof}

\begin{remark}[Connection to Linear Logic]
  Girard's linear logic~\cite{Girard} treats propositions as resources
  that are consumed in inference. The linear implication $A \multimap B$
  means ``given $A$, produce $B$, consuming $A$.'' This mirrors
  Spherepop events exactly: an event $\varepsilon : \Omega \to \Omega'$
  \emph{consumes} option space $\Omega$ (it can never happen again in
  the same form) and produces $\Omega'$. The pop operator is a linear
  implication. The technology tree is a derivation in linear logic.
\end{remark}

\begin{project}{Compositional World Transformations}
  \begin{enumerate}
    \item Implement the event composition operator
          $\varepsilon_2 \circ \varepsilon_1$ for pairs of world events.
    \item Verify associativity: confirm that three events applied in any
          associative grouping produce the same world state.
    \item Build a \texttt{VictoryChecker} and verify naturality
          (once true, stays true).
    \item Implement a \texttt{HistoryFunctor}: given an event log,
          compute the sequence of world states it produces.
  \end{enumerate}
\end{project}

% ── Chapter 28 ───────────────────────────────────────────────
\chapter{Sheaves and Local Coherence}
\label{ch:sheaves}

\begin{background}
  The world is large and agents within it are local. Each empire knows
  its territory and the rest poorly. Each city manages its own economy.
  Each unit sees only its vision radius. Yet the world must be globally
  coherent. The mathematical structure that captures this tension is a
  \emph{sheaf}~\cite{Spivak}: it assigns to each region an information
  structure such that sections on overlapping regions are consistent.
  Two overlapping regions produce the same information on their overlap.
  This is the mathematical model of distributed knowledge with
  local-global coherence.
\end{background}

\section{The Information Sheaf}

\begin{definition}[Information Sheaf]
  An \emph{information sheaf} $\mathcal{F}$ assigns to each region $U$
  a set $\mathcal{F}(U)$ of local sections (information knowable within $U$),
  and to each inclusion $V \subseteq U$ a restriction map
  $\rho_{UV} : \mathcal{F}(U) \to \mathcal{F}(V)$, satisfying:
  (1)~\textbf{Locality:} sections that agree on all cover elements are equal;
  (2)~\textbf{Gluing:} compatible local sections assemble into a unique
  global section.
\end{definition}

\begin{theorem}[Gluing Lemma as World Coherence]
  In a 4X world where each empire maintains local information, consistent
  local information can be assembled into a globally coherent world state
  if and only if the information structure forms a sheaf.
\end{theorem}
\begin{proof}
  The gluing axiom is precisely the condition needed: each empire's local
  information is internally consistent and agrees on shared borders.
  The gluing lemma provides a unique global section---a globally coherent
  world state---restricting to each empire's local information.
  Locality guarantees uniqueness. $\square$
\end{proof}

\begin{project}{Distributed AI Information Systems}
  \begin{enumerate}
    \item Implement a model where each empire stores state only for
          tiles in its explored region.
    \item Implement \texttt{BorderNegotiation}: reconcile conflicting
          information about border tiles between two empires.
    \item Implement \texttt{IntelligenceFusion}: combine spy reports
          to produce a coherent estimated world state.
    \item Build a \texttt{GlobalStateReconstructor}: assemble the full
          world state from all empires' local information.
  \end{enumerate}
\end{project}

% ── Chapter 29 ───────────────────────────────────────────────
\chapter{Semantic Infrastructure}
\label{ch:semantic}

\begin{background}
  The world simulation has become a \emph{semantic operating system}:
  a system in which the meaning of entities is determined by their
  relationships and histories, not by intrinsic properties. A city is
  meaningful not because of its data structure but because of its trade
  relationships, cultural influence, military history, and position in
  the civilization's narrative.

  This chapter makes the parallel explicit between Spherepop's event
  infrastructure and the event-sourced architectures of distributed
  computer systems: Git version control, Apache Kafka, distributed
  ledgers~\cite{Meijer,MeijerDuality}. Each solves the same fundamental
  problem: how to maintain coherent state across multiple participants
  who share a history but operate locally. The solution is always the
  same as in Spherepop: an append-only event log with defined merge and
  replay semantics.
\end{background}

\section{World Simulation as Distributed System}

\begin{definition}[Semantic Infrastructure]
  A \emph{semantic infrastructure} for world $\world$ is
  $\mathcal{SI} = (\mathcal{L}, \mathcal{R}, \mathcal{M}, \mathcal{A})$
  where $\mathcal{L}$ is the append-only event log, $\mathcal{R}$ is
  the replay function, $\mathcal{M}$ is the merge function, and
  $\mathcal{A}$ is the authority protocol for conflict resolution.
\end{definition}

\begin{theorem}[Semantic Infrastructure Completeness]
  A semantic infrastructure $\mathcal{SI}$ is \emph{complete} if for
  any two logs $\mathcal{L}_1, \mathcal{L}_2$ consistent with a common
  prefix $\mathcal{L}_0$, $\mathcal{M}(\mathcal{L}_1, \mathcal{L}_2)$
  is well-defined and contains $\mathcal{L}_0$ as a prefix.
\end{theorem}
\begin{proof}
  $\mathcal{M}(\mathcal{L}_1, \mathcal{L}_2)$ is defined as the
  interleaving of suffix events $\mathcal{L}_1 \setminus \mathcal{L}_0$
  and $\mathcal{L}_2 \setminus \mathcal{L}_0$, with conflicts resolved
  by $\mathcal{A}$. The result contains $\mathcal{L}_0$ as a prefix
  by construction. $\square$
\end{proof}

\begin{project}{Replayable Multiplayer World Architecture}
  \begin{enumerate}
    \item Implement a split-log architecture: each empire maintains its
          own local event log, synchronized with the global log each turn.
    \item Implement \texttt{LogMerge}: reconcile two empires' logs
          after simultaneous play.
    \item Implement the authority protocol: conflicts resolved by
          timestamp, then empire ID.
    \item Test coherence: two simultaneously-acting empires' logs merge
          into a consistent world state.
  \end{enumerate}
\end{project}

% ============================================================
% PART X: THE PHILOSOPHY OF COMPUTATIONAL WORLDS
% ============================================================

\part{The Philosophy of Computational Worlds}

\chapterquote{The limits of my language mean the limits of my world.
}{--- Wittgenstein~\cite{Wittgenstein}}

% ── Chapter 30 ───────────────────────────────────────────────
\chapter{Meaning as Use}
\label{ch:meaning}

\begin{background}
  Wittgenstein's later philosophy~\cite{Wittgenstein} proposes that the
  meaning of a word is its use in a language game---a structured social
  practice with rules, moves, and purposes. To understand ``castle''
  in chess, you do not need a theory of reference: you need to know
  how the castle moves. To understand ``city'' in a 4X game, you need
  to know what cities produce, what they cost, and how they interact
  with other entities.

  This chapter applies the Wittgensteinian framework to 4X world
  simulation. The meaning of in-game entities is constituted by the rule
  systems that govern them. Cultural evolution in the simulation is the
  evolution of language games. The diversity of civilizations corresponds
  to the diversity of mutually intelligible but distinct language games.
  A well-designed 4X world does not merely represent a civilization:
  it instantiates a form of life.
\end{background}

\section{Game Rules as Semantic Axioms}

\begin{definition}[Language Game]
  A \emph{language game} $\mathcal{G} = (\mathcal{M}, \mathcal{R},
  \mathcal{P})$ consists of moves $\mathcal{M}$, rules $\mathcal{R}$
  (constraints on legal move sequences), and purposes $\mathcal{P}$
  (goals that make sequences meaningful).
\end{definition}

\begin{proposition}[Game Rules as Semantics]
  The meaning of an in-game entity $E$ is the set of rules that govern
  $E$'s interactions: how it is produced, consumed, transformed, and
  used. Change the rules and you change the meaning, even if the data
  structure is identical.
\end{proposition}

\begin{mechanic}{Cultural Semantic Drift}
  As civilizations diverge in technology and diplomacy, their rule
  systems diverge. Two civilizations that cannot trade, communicate, or
  convert each other occupy different language games: they have lost the
  shared practices that would make mutual understanding possible.
  This is Wittgenstein's remark operationalized: ``If a lion could talk,
  we could not understand him.''
\end{mechanic}

\begin{project}{Cultural Semantics in Diplomacy}
  \begin{enumerate}
    \item Implement \texttt{DiplomaticIntelligibility} as a function of
          shared technologies, cultural traits, and religion.
    \item Build a rule that proposals require minimum intelligibility
          to be parseable by the target empire.
    \item Implement a \texttt{CulturalMissionary} unit that spreads
          traits, increasing intelligibility.
    \item Simulate two isolated civilizations developing for 40 turns
          then making first contact. Measure the intelligibility gap.
  \end{enumerate}
\end{project}

% ── Chapter 31 ───────────────────────────────────────────────
\chapter{Civilization as Computation}
\label{ch:civcomp}

\begin{background}
  We have arrived at the final synthesis. The entire book has been
  building toward a single claim: that a civilization is a computation.
  Not merely that a civilization can be \emph{modeled} by a
  computation---that is the weak claim of simulation---but that
  civilization \emph{is}, in its essential structure, a distributed
  information-processing system: a Turing-complete process that takes
  environmental inputs, maintains internal state, and produces outputs
  that modify its environment.

  This claim has a precise meaning and a precise proof. We will show
  that the 4X world simulation implemented in this book is Turing-complete,
  and therefore that civilization---the system it models---shares that
  property. The philosophical consequence is that civilization and
  computation are not analogous: they are the same thing at different
  scales of description, connected by the abstraction ladder we have
  climbed throughout this book.
\end{background}

\section{The Civilization as a Turing Machine}

\begin{theorem}[Civilization is Turing-Complete]
  The 4X world simulation implemented in this book is Turing-complete:
  it can simulate any Turing machine~\cite{Turing}.
\end{theorem}
\begin{proof}[Construction]
  We construct a Turing machine $\mathcal{TM}$ within the 4X world:
  \begin{itemize}
    \item \textbf{Tape:} A row of tiles; each tile's terrain type encodes
          a tape cell symbol $\in \{0, 1, \mathrm{blank}\}$.
    \item \textbf{Head:} A special unit (the ``reader unit'') whose
          position encodes the head position.
    \item \textbf{State register:} The reader unit's FSM state encodes
          the Turing machine's control state $q \in Q$.
    \item \textbf{Transition:} Each turn, the reader unit (i) reads the
          terrain type of its current tile; (ii) looks up the transition
          $(q, \mathrm{symbol}) \to (q', \mathrm{symbol}', \mathrm{dir})$
          from a table in the technology tree; (iii) writes the new symbol
          by changing terrain type; (iv) moves left or right by one tile.
  \end{itemize}
  This faithfully implements any Turing machine with tape alphabet
  $\{0, 1, \mathrm{blank}\}$ and finite state set $Q$. Since any Turing
  machine is simulated, the simulation is Turing-complete. $\square$
\end{proof}

\begin{perspective}
  Turing-completeness means that the internal logic of a civilization
  cannot, in general, be predicted without running the civilization.
  The future is genuinely open: there is no shortcut to history. This
  is Turing's halting problem applied to civilization. Whether an empire
  eventually achieves victory, collapses, or cycles indefinitely cannot
  be determined from initial conditions alone by any finite algorithm.
  The only way to know a world's history is to let it run.

  Wiener's cybernetics~\cite{Wiener} observed that control and
  communication in animals and machines obey the same principles.
  The Turing-completeness theorem is the strongest version of this
  claim: civilization and machine are not merely analogous but
  computationally equivalent.
\end{perspective}

\begin{project}{Self-Organizing Civilization Simulator}
  \begin{enumerate}
    \item Implement the Turing machine construction: tape tiles,
          reader unit, transition table.
    \item Program the reader unit to compute binary addition; verify
          the output.
    \item Implement \texttt{AutopoiesisCheck}: detect when an empire
          has reorganized itself without external input,
          following Varela et al.~\cite{Varela}.
    \item Run 10 simulations with random initial conditions; identify
          which exhibit sustained self-organization.
  \end{enumerate}
\end{project}

% ── Chapter 32 ───────────────────────────────────────────────
\chapter{The Finite World}
\label{ch:finite}

\chapterquote{The world is everything that is the case.
}{--- Wittgenstein, \textit{Tractatus}~\cite{Wittgenstein}}

\begin{background}
  Every world in this book is finite. The map has boundaries. The
  technology tree has leaves. The turn counter has a maximum. This
  finitude is not a deficiency---it is what makes the world computable.
  An infinite world cannot be serialized, replayed, or proven correct.
  Finitude is the price of intelligibility.

  But finitude also means that every world ends. Civilizations collapse
  or are conquered. Resources deplete. The final turn is reached. This
  last chapter examines what finitude means philosophically---for the
  simulation, for the civilizations within it, and for the beings outside
  it who have watched it run. It is the proper ending of a book that
  began with logic gates and arrives, finally, at the question of what
  all this construction means.
\end{background}

\section{Irreversibility and Path Dependence}

\begin{definition}[Path Dependence]
  A world $\world$ exhibits \emph{path dependence} at tick $t$ if there
  exist two event histories $\mathcal{L}_1 \neq \mathcal{L}_2$ consistent
  with the same initial conditions such that
  $\mathrm{state}(\world, \mathcal{L}_1, t) \neq
  \mathrm{state}(\world, \mathcal{L}_2, t)$.
\end{definition}

\begin{theorem}[All Non-Trivial Worlds are Path Dependent]
  Any world with at least one irreversible event exhibits path dependence
  at all subsequent ticks.
\end{theorem}
\begin{proof}
  Let $\varepsilon$ be an irreversible event. Let
  $\mathcal{L}_1 = (\varepsilon)$ and $\mathcal{L}_2 = ()$ (empty log).
  By irreversibility, $\mathrm{state}(\mathcal{L}_1) \neq
  \mathrm{state}(\mathcal{L}_2)$ at tick $t \geq t_\varepsilon$.
  $\square$
\end{proof}

\section{Historical Lock-In}

\begin{definition}[Historical Lock-In]
  A civilization is \emph{locked in} at tick $t$ if
  $\Phi_C(t) \leq \theta_{\min}$ for some minimum viable threshold
  $\theta_{\min}$: all strategic variation has been exhausted.
\end{definition}

\begin{proposition}[Lock-In is Generically Monotone]
  A civilization that has lost all strategic flexibility cannot, by
  definition, take the actions that would restore flexibility.
\end{proposition}

\section{The Final Perspective}

\begin{perspective}
  This book began with a logic gate. It ends with a civilization.

  Between those two points lies a philosophical ascent that mirrors, in
  miniature, the ascent of complexity in the actual universe: from the
  simplest distinctions through the accumulation of memory, the emergence
  of space, the arising of agency, the construction of economies and
  treaties and histories, to the moment when the whole system becomes
  so intricate that its future cannot be predicted from its present.

  What have we built? Not merely a game. We have built a system in which
  local rules, applied consistently across space and time, produce global
  properties that no individual rule anticipated: Zipf-distributed city
  sizes, percolation thresholds in territorial connectivity, diplomatic
  trust dynamics resembling reputation equilibria, cultural drift
  resembling linguistic evolution. We have built a small model of
  history---not history with specific content, but history as a
  \emph{form}: the form of irreversible, accumulating, locally-caused,
  globally-surprising change.

  Heidegger~\cite{Heidegger} observed that the question of Being is the
  question every other question presupposes. For computational worlds,
  the question of Being is the question we have answered: what does it
  mean for an entity to \emph{be} in a computational world? It means:
  to have a history. To have been caused. To be maintained by the
  persistence of its effects. To be, in the Spherepop sense, an
  append-only log that has not yet been closed.

  The world runs. The log grows. The future remains, for now, open.
\end{perspective}

\begin{project}{Complete Historical Computational Civilization}
  \textbf{Grand Final Project.}

  Build a complete, persistent, historically self-aware 4X civilization
  simulator incorporating every system from Chapters 1--31. Requirements:
  \begin{enumerate}
    \item A $32 \times 32$ world with procedural terrain and ecology.
    \item Four empires with distinct cultural trait vectors and
          starting technologies.
    \item Full event sourcing: all state derived from event logs.
    \item Optionality tracking for all empires at every turn.
    \item A narrative generator producing a structured chronicle.
    \item A save/load/replay system with full integrity verification.
    \item A post-game analysis: which decisions caused the largest
          option-space contractions? Which empire maintained optionality
          longest? What was the Zipf exponent of the final city
          distribution? When did each empire approach lock-in?
  \end{enumerate}

  \medskip
  \noindent\textit{A student who completes this project has not merely
  learned to program a game. They have built a machine for generating
  history, and in doing so, have understood something about what
  history is.}
\end{project}

% ============================================================
% APPENDICES A--H (with full proofs)
% ============================================================



% ============================================================
% APPENDICES
% ============================================================

\appendix


\chapter{Spherepop Language Reference}
\label{app:spherepop}

This appendix provides a complete reference for the \Sp{} notation used
throughout the book.

\section{Primitive Bubbles}

\begin{center}
\begin{tabular}{lll}
  \toprule
  \textbf{Name} & \textbf{Signature} & \textbf{Semantics} \\
  \midrule
  \texttt{NAND} & $(a, b) \to \texttt{out}$
                & $\texttt{out} = \NOT(a \AND b)$ \\
  \texttt{DFF}  & $(\texttt{in}) \to \texttt{out}$
                & $\texttt{out}(t) = \texttt{in}(t-1)$ \\
  \texttt{MUX}  & $(a, b, \texttt{sel}) \to \texttt{out}$
                & $\texttt{out} = a$ if $\texttt{sel}=0$ else $b$ \\
  \texttt{DMUX} & $(\texttt{in}, \texttt{sel}) \to (a, b)$
                & $(a,b) = (\texttt{in},0)$ if $\texttt{sel}=0$ else $(0,\texttt{in})$ \\
  \bottomrule
\end{tabular}
\end{center}

\section{Syntax Summary}

\begin{lstlisting}[caption={Bubble declaration syntax}]
bubble BubbleName {
    -- Port declarations
    inputs:  portA, portB[n]   -- scalar or n-bit bus
    output:  portC, portD[m]

    -- State declarations (persistent across ticks)
    latch name[width]

    -- Combinational logic
    name = expression

    -- Sequential logic (executes at tick boundary)
    when condition do
        name := expression    -- := denotes sequential assignment
    end

    -- Iteration
    let i = 0
    while i < n do
        ...
        i = i + 1
    end

    -- Sub-bubble instantiation
    let result = SubBubble(arg1, arg2)
}
\end{lstlisting}

\section{Built-in Expressions}

\begin{center}
\begin{tabular}{ll}
  \toprule
  \textbf{Expression} & \textbf{Meaning} \\
  \midrule
  \texttt{a and b}  & Boolean AND \\
  \texttt{a or b}   & Boolean OR \\
  \texttt{not a}    & Boolean NOT \\
  \texttt{a + b}    & Integer addition (modular) \\
  \texttt{a - b}    & Integer subtraction (clamped at 0) \\
  \texttt{a >= b}   & Comparison (returns Boolean) \\
  \texttt{a[k]}     & Bit $k$ of bus $a$ \\
  \texttt{a[j..k]}  & Bit slice \\
  \bottomrule
\end{tabular}
\end{center}

\chapter{Truth Tables and Gate Diagrams}
\label{app:truthtables}

\section{Fundamental Gates}

\begin{multicols}{2}

\noindent\textbf{AND}\\[4pt]
\begin{tabular}{cc|c}
  \toprule
  $a$ & $b$ & \texttt{out} \\
  \midrule
  0 & 0 & 0 \\
  0 & 1 & 0 \\
  1 & 0 & 0 \\
  1 & 1 & 1 \\
  \bottomrule
\end{tabular}

\columnbreak

\noindent\textbf{OR}\\[4pt]
\begin{tabular}{cc|c}
  \toprule
  $a$ & $b$ & \texttt{out} \\
  \midrule
  0 & 0 & 0 \\
  0 & 1 & 1 \\
  1 & 0 & 1 \\
  1 & 1 & 1 \\
  \bottomrule
\end{tabular}

\end{multicols}

\vspace{1em}

\begin{multicols}{2}

\noindent\textbf{NAND}\\[4pt]
\begin{tabular}{cc|c}
  \toprule
  $a$ & $b$ & \texttt{out} \\
  \midrule
  0 & 0 & 1 \\
  0 & 1 & 1 \\
  1 & 0 & 1 \\
  1 & 1 & 0 \\
  \bottomrule
\end{tabular}

\columnbreak

\noindent\textbf{XOR}\\[4pt]
\begin{tabular}{cc|c}
  \toprule
  $a$ & $b$ & \texttt{out} \\
  \midrule
  0 & 0 & 0 \\
  0 & 1 & 1 \\
  1 & 0 & 1 \\
  1 & 1 & 0 \\
  \bottomrule
\end{tabular}

\end{multicols}

\chapter{The 4X Design Pattern Glossary}
\label{app:glossary}

\begin{description}[style=nextline, leftmargin=2cm]
  \item[Agent]
    A bubble with internal state, perception, and action capability.
    The computational model of a unit, city, or empire.

  \item[Boolean permission]
    A conjunction of conditions that must all hold for an action to
    be legal. The fundamental structure of game rules.

  \item[Bubble]
    The basic computational unit of \Sp{}: a bounded, ported,
    composable entity that simultaneously models data, process, and
    environment.

  \item[Cellular automaton]
    A spatial update system where each cell updates based on its
    local neighborhood. The computational model of physical laws
    operating on the world map.

  \item[Event log]
    An append-only trace of world state transitions. The computational
    model of history.

  \item[Fog of war]
    The visibility restriction that limits an agent's perception to
    explored and currently visible tiles. Implemented as a pair of
    Boolean registers per tile.

  \item[Latch / Register]
    A persistent storage element, built from feedback circuits.
    The fundamental implementation of world memory.

  \item[Technology DAG]
    A directed acyclic graph of Boolean predicates representing
    technological dependencies. Prerequisite checking is conjunction
    over parent predicates.

  \item[Trade route]
    A persistent path signal between two city bubbles. Disruption
    is implemented as path invalidation.

  \item[Treaty]
    A shared mutable bubble between two empire agents, storing
    diplomatic state.

  \item[Victory condition]
    A global Boolean predicate over the entire world state.
    The terminal condition of a simulation run.
\end{description}




\chapter{Formal Bubble Semantics}
\label{app:bubble-semantics}

\chapterquote{%
  A bubble is a stateful morphism: it transforms inputs into outputs
  while maintaining internal memory.
}{}

The preceding appendices treated bubbles operationally. This appendix
gives the algebraic definition that grounds the entire \Sp{} architecture.

\section{Bubbles as Mealy Machines}

\begin{definition}[Bubble]
  A \emph{Spherepop bubble} is a tuple
  \[
    B = (I, O, S, \delta, \lambda, s_0)
  \]
  where:
  \begin{itemize}
    \item $I = \mathbb{B}^m$ is the \emph{input port space} ($m$-bit input vector),
    \item $O = \mathbb{B}^p$ is the \emph{output port space} ($p$-bit output vector),
    \item $S = \mathbb{B}^q$ is the \emph{internal state space} ($q$-bit state register),
    \item $\delta : S \times I \to S$ is the \emph{state transition function}
          (next-state logic),
    \item $\lambda : S \times I \to O$ is the \emph{output function}
          (output logic),
    \item $s_0 \in S$ is the \emph{initial state}.
  \end{itemize}
\end{definition}

This is precisely the definition of a \emph{Mealy machine} (finite-state
transducer). Mealy machines are the standard model for synchronous
sequential digital circuits. Every bubble in \Sp{} is a Mealy machine;
the entire world is the product of all its bubbles' Mealy machines,
synchronized by the turn counter.

\begin{proposition}[Combinational Gates are Bubbles]
  A combinational gate (no internal state) is a bubble with $S = \{s_0\}$
  (singleton state), $\delta(s_0, i) = s_0$ for all $i$, and
  $\lambda(s_0, i) = f(i)$ where $f : I \to O$ is the gate's Boolean function.
\end{proposition}

\begin{proposition}[Sequential Memory is a Bubble]
  A 1-bit D flip-flop is the bubble $(I, O, \mathbb{B}, \delta, \lambda, 0)$
  with $I = O = \mathbb{B}$, $\delta(s, i) = i$, $\lambda(s, i) = s$.
  The output is the \emph{previous tick's} input, implementing a 1-tick delay.
\end{proposition}

\section{Bubble Composition}

\begin{definition}[Sequential Composition]
  Given bubbles $B_1 = (I_1, O_1, S_1, \delta_1, \lambda_1, s_1^0)$ and
  $B_2 = (I_2, O_2, S_2, \delta_2, \lambda_2, s_2^0)$ with $O_1 = I_2$,
  their \emph{sequential composition} $B_2 \circ B_1$ is the bubble
  \[
    (I_1,\; O_2,\; S_1 \times S_2,\; \delta,\; \lambda,\; (s_1^0, s_2^0))
  \]
  where:
  \begin{align*}
    \delta((s_1, s_2), i) &= (\delta_1(s_1, i),\; \delta_2(s_2, \lambda_1(s_1, i))) \\
    \lambda((s_1, s_2), i) &= \lambda_2(s_2, \lambda_1(s_1, i))
  \end{align*}
\end{definition}

\begin{theorem}[Composition Associativity]
  Sequential bubble composition is associative:
  $(B_3 \circ B_2) \circ B_1 = B_3 \circ (B_2 \circ B_1)$.
\end{theorem}
\begin{proof}
  Both sides implement the same input-output function: apply $B_1$'s
  output as $B_2$'s input, then $B_2$'s output as $B_3$'s input.
  State spaces are isomorphic by associativity of Cartesian product.
  $\square$
\end{proof}

\section{The Category of Bubbles}

\begin{definition}[Bubble Category $\mathbf{Bub}$]
  The category $\mathbf{Bub}$ has:
  \begin{itemize}
    \item \textbf{Objects:} port types $\mathbb{B}^n$ for $n \geq 0$.
    \item \textbf{Morphisms:} $\mathrm{Hom}(\mathbb{B}^m, \mathbb{B}^p)$
          is the set of bubbles with input space $\mathbb{B}^m$ and
          output space $\mathbb{B}^p$.
    \item \textbf{Composition:} sequential bubble composition.
    \item \textbf{Identity:} for each $\mathbb{B}^n$, the identity bubble
          with $\delta = \mathrm{id}$ and $\lambda(s,i) = i$.
  \end{itemize}
\end{definition}

\begin{proposition}[$\mathbf{Bub}$ is a Category]
  $\mathbf{Bub}$ satisfies the category axioms: composition is associative
  (by Theorem above) and identities exist.
\end{proposition}

\begin{remark}
  The world evolution operator $\Fworld$ (Chapter~\ref{sec:worldtransition})
  is an endomorphism in $\mathbf{Bub}$:
  $\Fworld \in \mathrm{Hom}(\mathbb{B}^N, \mathbb{B}^N)$ where $N$ is the
  total number of state bits in the world. The world's entire history is
  the orbit $\{s_0, \Fworld(s_0), \Fworld^2(s_0), \ldots\}$ under
  repeated application of this morphism.
\end{remark}

\section{Connection to Markov Blankets}

The Mealy machine formulation makes the Markov blanket structure
(Chapter~\ref{ch:blankets}) precise:

\begin{proposition}[Bubble Blanket Correspondence]
  For bubble $B = (I, O, S, \delta, \lambda, s_0)$:
  \begin{align*}
    \mathcal{I} &= S \quad \text{(internal state registers)} \\
    \mathcal{S} &= I \quad \text{(sensory states = input ports)} \\
    \mathcal{A} &= O \quad \text{(active states = output ports)} \\
    \mathcal{B} &= I \cup O \quad \text{(the blanket = all ports)} \\
    \mathcal{E} &= \text{all other bubble states in } \world
  \end{align*}
  The state transition $\delta : S \times I \to S$ updates internal
  states given blanket input (active inference). The output function
  $\lambda : S \times I \to O$ generates active states from internal
  states (acting on the world).
\end{proposition}

This shows that the bubble definition already encodes the full
active inference loop: internal states infer from sensory inputs,
and active outputs modify the external world.


\chapter{Boolean Foundations of Spherepop Gates}
\label{app:boolean}

\section{Primitive Binary Algebra}

Let $\mathbb{B} = \{0, 1\}$. A Spherepop primitive gate is a morphism
$g : \mathbb{B}^n \to \mathbb{B}^m$.

\begin{definition}[Gate Graph]
  A \emph{Spherepop gate graph} is a finite directed acyclic graph
  $G = (V, E, \lambda)$ where $V$ is the set of gate instances,
  $E$ is the set of signal edges, and $\lambda : V \to \mathcal{G}$
  assigns a primitive gate type to each node.
\end{definition}

\begin{theorem}[Functional Completeness of NAND]
  Every Boolean function $f : \mathbb{B}^n \to \mathbb{B}^m$ is
  expressible as a finite composition of NAND gates.
\end{theorem}
\begin{proof}
  It suffices to derive NOT and AND from NAND (since NOT and AND are
  functionally complete).
  \begin{align*}
    \mathrm{NOT}(a) &= \mathrm{NAND}(a, a). \\
    &\textit{Check: } \mathrm{NAND}(0,0)=1=\mathrm{NOT}(0);\;
      \mathrm{NAND}(1,1)=0=\mathrm{NOT}(1). \\
    \mathrm{AND}(a,b) &= \mathrm{NOT}(\mathrm{NAND}(a,b))
                        = \mathrm{NAND}(\mathrm{NAND}(a,b),\,\mathrm{NAND}(a,b)). \\
    &\textit{Check: } \mathrm{AND}(1,1)=\mathrm{NAND}(\mathrm{NAND}(1,1),\mathrm{NAND}(1,1))
      =\mathrm{NAND}(0,0)=1. \checkmark
  \end{align*}
  OR follows via De Morgan: $a \vee b = \neg(\neg a \wedge \neg b)$.
  Since $\{\neg, \wedge\}$ is functionally complete, NAND alone suffices.
  $\square$
\end{proof}

\section{Evaluation Semantics}

\begin{theorem}[Evaluation is Well-Defined]
  For any acyclic gate graph $G$, topological evaluation produces a
  unique output assignment, independent of the choice of topological
  ordering.
\end{theorem}
\begin{proof}
  A topological ordering exists (since $G$ is acyclic). For any two
  topological orderings, each gate is evaluated after all its predecessors
  in both orderings, receiving the same input signals in both, and
  therefore producing the same output. Uniqueness follows from the
  determinism of each primitive gate. $\square$
\end{proof}

% ── Appendix B ───────────────────────────────────────────────
\chapter{Event Histories and Option Spaces}
\label{app:events}

\section{Historical State Spaces}

\begin{definition}[Option Space Dynamics]
  A world history is the compositional chain
  $H = \varepsilon_n \circ \cdots \circ \varepsilon_1$
  with monotone admissibility: $X_{t+1} \subseteq X_t$.
\end{definition}

\begin{theorem}[Monotone Option Space Reduction]
  For any non-trivial event sequence of length $n$:
  $|X_n| \leq |X_0|$, with strict inequality if at least one event
  is non-trivial.
\end{theorem}
\begin{proof}
  Induction on $n$. Base: $|X_0| \leq |X_0|$. Inductive step: assume
  $|X_k| \leq |X_0|$. By monotone admissibility,
  $|X_{k+1}| \leq |X_k| \leq |X_0|$.
  Strict inequality follows from Strict Contraction (Chapter~21)
  for any non-trivial $\varepsilon_{k+1}$. $\square$
\end{proof}

\section{Historical Identity as Equivalence Relation}

\begin{proposition}
  Historical identity $\equiv_H$ is an equivalence relation.
\end{proposition}
\begin{proof}
  Reflexivity: $\mathcal{L}(A) = \mathcal{L}(A)$, so $A \equiv_H A$.
  Symmetry: $\mathcal{L}(A)=\mathcal{L}(B) \Rightarrow
  \mathcal{L}(B)=\mathcal{L}(A)$.
  Transitivity: $\mathcal{L}(A)=\mathcal{L}(B)=\mathcal{L}(C)
  \Rightarrow A \equiv_H C$. $\square$
\end{proof}

% ── Appendix C ───────────────────────────────────────────────
\chapter{Merge--Collapse Algebra}
\label{app:merge}

\section{Merge Lattice}

\begin{theorem}[Territory Merge Lattice]
  $(\mathcal{P}(\mathrm{Tiles}), \cup, \cap)$ is a bounded
  distributive lattice.
\end{theorem}
\begin{proof}
  Commutativity, associativity, and idempotence of $\cup$ and $\cap$
  are standard. Absorption: $A \cup (A \cap B) = A$ and
  $A \cap (A \cup B) = A$. Distributivity:
  $A \cup (B \cap C) = (A \cup B) \cap (A \cup C)$---a standard
  set identity. Bottom: $\emptyset$; top: $\mathrm{Tiles}$. $\square$
\end{proof}

\begin{theorem}[Collapse Normalization]
  The collapse operator $\kappa$ is a lattice homomorphism:
  $\kappa(\mu(A,B)) = \mu(\kappa(A), \kappa(B))$.
\end{theorem}
\begin{proof}
  $\kappa$ maps each element to its equivalence class representative.
  Since $\mu$ is defined on equivalence classes (merging representatives
  gives the representative of the merged class), $\kappa$ commutes
  with $\mu$. $\square$
\end{proof}

% ── Appendix D ───────────────────────────────────────────────
\chapter{Sequential Bubble Machines}
\label{app:sequential}

\section{Bubble Automata}

\begin{definition}[History-Sensitive Automaton]
  A sequential Spherepop machine is
  $\mathcal{M} = (Q, \Sigma, \Gamma, \delta, q_0, F, H)$
  where $\delta : Q \times \Sigma \times H \to Q \times \Gamma \times H$
  is history-sensitive (identical local states may produce different
  futures depending on prior commitments in $H$).
\end{definition}

\begin{theorem}[History-Sensitivity Strictly Increases Power]
  History-sensitive automata are strictly more expressive than standard
  finite automata.
\end{theorem}
\begin{proof}
  \textit{Containment:} Any FA is simulated by ignoring $H$.
  \textit{Strictness:} Consider
  $L = \{w \in \{a,b\}^* : \text{first } a \text{ was preceded by an
  even number of } b\text{s}\}$.
  No FA accepts $L$ (pumping lemma: the count of $b$s before the first
  $a$ is unbounded). A history-sensitive automaton accepts $L$ by
  scanning $H$ to count $b$-events before the first $a$-event.
  $\square$
\end{proof}

% ── Appendix E ───────────────────────────────────────────────
\chapter{Entropy and Optionality}
\label{app:entropy}

\section{Historical Entropy}

\begin{theorem}[Pop Operator Satisfies Second Law]
  For any non-trivial pop $\mathbf{P}_\varepsilon$:
  $S(x_{t+1}) \leq S(x_t)$.
\end{theorem}
\begin{proof}
  $|\mathcal{T}(x_{t+1})| = |\mathcal{T}(x_t) \cap \Omega'|
  \leq |\mathcal{T}(x_t)|$. Since $\log_2$ is monotone,
  $S(x_{t+1}) \leq S(x_t)$. $\square$
\end{proof}

This is the formal analogue of Landauer's principle~\cite{Landauer}:
logical irreversibility corresponds to thermodynamic entropy production.

\section{Collapse Criterion}

\begin{theorem}[Collapse Criterion]
  Civilization $C$ is in irreversible collapse at tick $t$ if
  $\Phi_C(t) < \Phi_{\min}$ and $\Delta\Phi_C(t) < 0$ simultaneously.
  Under these conditions $\Phi_C(t') < \Phi_{\min}$ for all $t' > t$
  in expectation.
\end{theorem}
\begin{proof}
  If $\Phi_C < \Phi_{\min}$, the civilization lacks the option diversity
  to mount effective responses. Each turn, forced pops (responses to
  attacks, famines) reduce $\Phi_C$ without the civilization being able
  to choose high-value expansion pops. Thus $\Delta\Phi_C < 0$ compounds
  until $\Phi_C \to 0$. $\square$
\end{proof}

% ── Appendix F ───────────────────────────────────────────────
\chapter{Category-Theoretic Semantics}
\label{app:category}

\section{The History Category}

\begin{theorem}[$\mathbf{Hist}$ is a Monoid Category]
  The endomorphisms of a fixed option space $\Omega$:
  $\mathrm{End}(\Omega) = \{\varepsilon : \Omega \to \Omega\}$
  form a monoid under composition with identity $\mathrm{id}_\Omega$.
\end{theorem}
\begin{proof}
  Closure: composition of two maps $\Omega \to \Omega$ gives
  $\Omega \to \Omega$. Associativity: function composition is
  associative. Identity: $\mathrm{id}_\Omega$ is the unit. $\square$
\end{proof}

\begin{theorem}[No Traced Monoidal Structure Without Path Erasure]
  $\mathbf{Hist}$ does not admit a trace operator~\cite{JoyalStreetVerity}
  consistent with the irreversibility constraint.
\end{theorem}
\begin{proof}
  A trace operator would allow feedback: a future event causing a past
  event to be undone. The trace of a non-trivial pop
  $\varepsilon : \Omega \otimes U \to \Omega' \otimes U$ with
  $\Omega' \subsetneq \Omega$ would require the feedback wire to
  ``undo'' the option-space contraction, contradicting Strict Contraction.
  $\square$
\end{proof}

\section{Monotone Victory Conditions}

\begin{proposition}
  Monotone victory conditions form a subposet of the lattice of natural
  transformations $F \Rightarrow \mathbf{2}$.
\end{proposition}
\begin{proof}
  A natural transformation $V : F \Rightarrow \mathbf{2}$ is monotone
  if $V_\Omega(s)=1 \Rightarrow V_{\Omega'}(F(\varepsilon)(s))=1$
  for all events $\varepsilon$. The set of monotone natural
  transformations is closed under pointwise meet and join, forming
  a subposet. $\square$
\end{proof}

% ── Appendix G ───────────────────────────────────────────────
\chapter{Procedural World Dynamics}
\label{app:pde}

\section{Tile Evolution PDE}

\begin{definition}[World Evolution Equation]
  Resource density $\rho(x,y,t)$ evolves as:
  \[
    \frac{\partial \rho}{\partial t}
    = D \nabla^2 \rho + r\rho\!\left(1 - \frac{\rho}{K}\right)
    - c \cdot A \cdot \rho
  \]
  where $D$ is the diffusion coefficient, $r$ is ecological growth rate,
  $K$ is carrying capacity, $c$ is per-capita extraction rate, and $A$
  is agent population density.
\end{definition}

\begin{theorem}[Discrete Stability Condition]
  The forward-Euler discretization is stable if
  $\Delta t \leq h^2 / (4D)$ (2D CFL condition).
\end{theorem}
\begin{proof}
  The 2D diffusion stencil gives amplification factor
  $|1 - 4D\Delta t/h^2|$. Stability requires this $\leq 1$, giving
  $\Delta t \leq h^2/(2D)$ in 1D and $\Delta t \leq h^2/(4D)$ in 2D
  (the factor of 2 from the two spatial dimensions). $\square$
\end{proof}

\begin{theorem}[Equilibrium Resource Density]
  With $A=0$, resource density converges to $\rho^* = K$ if $\rho_0 > 0$.
\end{theorem}
\begin{proof}
  With $D=0$, each tile obeys $d\rho/dt = r\rho(1-\rho/K)$, the logistic
  equation with stable equilibrium $\rho^* = K$. With $D > 0$, diffusion
  averages spatial heterogeneity, driving the system toward the uniform
  equilibrium exponentially. $\square$
\end{proof}

% ── Appendix H ───────────────────────────────────────────────
\chapter{Trajectory Collapse Geometry}
\label{app:trajectory}

\section{Hypothesis Manifolds}

\begin{theorem}[Maximum Entropy Prior]
  In the absence of distinguishing information, the maximum entropy
  distribution on $\Delta_{n-1}$ is uniform: $p_i = 1/n$, $H = \log_2 n$.
\end{theorem}
\begin{proof}
  By the Gibbs inequality, $H(\mathbf{q}) \leq \log_2 n$ for any
  distribution on $n$ elements, with equality iff $q_i = 1/n$.
  $\square$
\end{proof}

\section{Information Gain from Strategic Intelligence}

\begin{definition}[Strategic Information Gain]
  An intelligence report $r$ provides information gain:
  \[
    \mathrm{IG}(r) = H(\mathbf{p}) - H(\mathbf{p}' \mid r).
  \]
\end{definition}

\begin{corollary}
  Perfect intelligence about scenario $h_k$ (setting $p_k' = 1$)
  provides maximum information gain $\mathrm{IG}_{\max} = \log_2 n$ bits.
\end{corollary}

This formalizes the intuition that espionage is valuable precisely to
the extent that it resolves strategic uncertainty. An empire that spends
on intelligence is buying bits of future option-space clarity.


\backmatter

\chapter*{Bibliography}
\addcontentsline{toc}{chapter}{Bibliography}

\begin{thebibliography}{99}

\bibitem{Aristotle}
Aristotle.
\textit{Nicomachean Ethics}.
Translated by W.~D.~Ross.
Oxford University Press, 2009.

\bibitem{Badiou}
Alain Badiou.
\textit{Being and Event}.
Continuum, 2005.

\bibitem{Barendregt}
Henk Barendregt.
\textit{The Lambda Calculus: Its Syntax and Semantics}.
North-Holland, 1984.

\bibitem{Bartosz}
Bartosz Milewski.
\textit{Category Theory for Programmers}.
Blurb, 2018.

\bibitem{BaezStay}
John~C.\ Baez and Mike Stay.
Physics, topology, logic and computation: A Rosetta Stone.
In \textit{New Structures for Physics}, pages 95--172.
Springer, 2010.

\bibitem{Church}
Alonzo Church.
An unsolvable problem of elementary number theory.
\textit{American Journal of Mathematics},
58(2):345--363, 1936.

\bibitem{Fine}
Kit Fine.
Things and their parts.
\textit{Midwest Studies in Philosophy},
23(1):61--74, 1999.

\bibitem{Foucault}
Michel Foucault.
\textit{The Archaeology of Knowledge}.
Pantheon Books, 1972.

\bibitem{Girard}
Jean-Yves Girard.
Linear logic.
\textit{Theoretical Computer Science},
50(1):1--101, 1987.

\bibitem{Heidegger}
Martin Heidegger.
\textit{Being and Time}.
Harper Perennial, 2008.

\bibitem{JoyalStreetVerity}
Andr\'{e} Joyal, Ross Street, and Dominic Verity.
Traced monoidal categories.
\textit{Mathematical Proceedings of the Cambridge Philosophical Society},
119(3):447--468, 1996.

\bibitem{Kim}
Jaegwon Kim.
Events as property exemplifications.
In \textit{Supervenience and Mind}.
Cambridge University Press, 1993.

\bibitem{Landauer}
Rolf Landauer.
Irreversibility and heat generation in the computing process.
\textit{IBM Journal of Research and Development},
5(3):183--191, 1961.

\bibitem{MacLane}
Saunders Mac~Lane.
\textit{Categories for the Working Mathematician}.
Springer, 1998.

\bibitem{Meijer}
Erik Meijer.
Your mouse is a database.
Keynote presentation, Build Conference, Microsoft, 2014.

\bibitem{MeijerDuality}
Erik Meijer and Gavin Bierman.
A co-relational model of data for large shared data banks.
Lectures on duality, category theory, and reactive computation, 2011--2015.

\bibitem{Milner}
Robin Milner.
\textit{Communicating and Mobile Systems: The Pi-Calculus}.
Cambridge University Press, 1999.

\bibitem{NisanSchocken}
Noam Nisan and Shimon Schocken.
\textit{The Elements of Computing Systems:
Building a Modern Computer from First Principles}.
MIT Press, 2005.

\bibitem{Shannon}
Claude~E.\ Shannon.
A mathematical theory of communication.
\textit{Bell System Technical Journal},
27(3):379--423, 1948.

\bibitem{Spivak}
David~I.\ Spivak.
\textit{Category Theory for the Sciences}.
MIT Press, 2014.

\bibitem{Turing}
Alan~M.\ Turing.
On computable numbers, with an application to the Entscheidungsproblem.
\textit{Proceedings of the London Mathematical Society},
42(2):230--265, 1936.

\bibitem{Varela}
Francisco Varela, Humberto Maturana, and Ricardo Uribe.
Autopoiesis: The organization of living systems.
\textit{BioSystems},
5(4):187--196, 1974.

\bibitem{VonNeumann}
John von Neumann.
\textit{Theory of Self-Reproducing Automata}.
University of Illinois Press, 1966.

\bibitem{Whitehead}
Alfred North Whitehead.
\textit{Process and Reality}.
Free Press, 1978.

\bibitem{Wiener}
Norbert Wiener.
\textit{Cybernetics: Or Control and Communication in the Animal and the Machine}.
MIT Press, 1948.

\bibitem{Wittgenstein}
Ludwig Wittgenstein.
\textit{Philosophical Investigations}.
Blackwell Publishing, 1953.

\bibitem{Wolfram}
Stephen Wolfram.
\textit{A New Kind of Science}.
Wolfram Media, 2002.

\bibitem{Wilson}
Mark Wilson.
\textit{Wandering Significance: An Essay on Conceptual Behaviour}.
Oxford University Press, 2006.

\bibitem{Axelrod}
Robert Axelrod.
\textit{The Evolution of Cooperation}.
Basic Books, 1984.

\bibitem{Friston}
Karl Friston.
The free-energy principle: a unified brain theory?
\textit{Nature Reviews Neuroscience},
11(2):127--138, 2010.

\bibitem{FristonBlanket}
Karl Friston.
Life as we know it.
\textit{Journal of the Royal Society Interface},
10(86):20130475, 2013.

\bibitem{Levin}
Michael Levin.
Technological approach to mind everywhere:
an experimentally grounded framework for understanding
diverse bodies and minds.
\textit{Frontiers in Systems Neuroscience},
16:768201, 2022.

\end{thebibliography}

\chapter*{Index of Spherepop Bubbles}
\addcontentsline{toc}{chapter}{Index of Spherepop Bubbles}

The following bubbles are defined or specified in this textbook:

\begin{multicols}{2}
\begin{itemize}[leftmargin=*, itemsep=0pt]
  \item \texttt{NAND} --- Ch.~1
  \item \texttt{NOT}, \texttt{AND}, \texttt{OR} --- Ch.~1
  \item \texttt{XOR}, \texttt{MUX}, \texttt{DMUX} --- Ch.~1
  \item \texttt{PermissionCheck} --- Ch.~1
  \item \texttt{AND16}, \texttt{OR16}, \texttt{NOT16} --- Ch.~2
  \item \texttt{Add16}, \texttt{Inc16} --- Ch.~2
  \item \texttt{ResourceAdder} --- Ch.~2
  \item \texttt{PopulationALU} --- Ch.~2
  \item \texttt{Bit}, \texttt{Register16} --- Ch.~3
  \item \texttt{RAM8} \ldots \texttt{RAM16K} --- Ch.~3
  \item \texttt{TileState}, \texttt{WorldMap} --- Ch.~3, 5
  \item \texttt{CellularAutomaton} --- Ch.~4
  \item \texttt{TerrainDecoder}, \texttt{TileYield} --- Ch.~5
  \item \texttt{FoodIntegrator}, \texttt{GoldLedger} --- Ch.~6
  \item \texttt{ResearchAccumulator} --- Ch.~6
  \item \texttt{CityProduction} --- Ch.~6
  \item \texttt{ChebyshevCheck} --- Ch.~7
  \item \texttt{VisionBroadcaster}, \texttt{FogUpdater} --- Ch.~7
  \item \texttt{ExplorationTracker} --- Ch.~7
  \item \texttt{Unit}, \texttt{UnitFSM} --- Ch.~8
  \item \texttt{MoveValidator}, \texttt{CombatResolver} --- Ch.~8
  \item \texttt{SettlePermission}, \texttt{CityFounder} --- Ch.~9
  \item \texttt{BorderExpander} --- Ch.~9
  \item \texttt{PathChecker}, \texttt{TradeRoute} --- Ch.~10
  \item \texttt{GoldCalculator} --- Ch.~10
  \item \texttt{DamageCalculator}, \texttt{CombatRound} --- Ch.~11
  \item \texttt{CaptureCity} --- Ch.~11
  \item \texttt{Treaty}, \texttt{TrustMatrix} --- Ch.~12
  \item \texttt{DiplomacyAI} --- Ch.~12
  \item \texttt{TechTree}, \texttt{PrereqChecker} --- Ch.~13
  \item \texttt{ForestSpread}, \texttt{DepletionTracker} --- Ch.~14
  \item \texttt{ClimateLayer} --- Ch.~14
  \item \texttt{EventLog}, \texttt{HistoryQuery} --- Ch.~15
  \item \texttt{NarrativeGenerator} --- Ch.~15
  \item \texttt{TurnManager}, \texttt{WorldBubble} --- Ch.~16
  \item \texttt{SaveSystem} --- Ch.~17
  \item \texttt{CheckpointManager} --- Ch.~17
  \item \texttt{EventSourcedWorld} --- Ch.~19
  \item \texttt{QueryOwner}, \texttt{AppendEvent} --- Ch.~19
  \item \texttt{TrustFromHistory} --- Ch.~19
  \item \texttt{TimelineReconstructor} --- Ch.~23
  \item \texttt{CityBiography} --- Ch.~23
  \item \texttt{ConnectedComponents} --- Ch.~24
  \item \texttt{CutVertexDetector} --- Ch.~24
  \item \texttt{Compactness} --- Ch.~24
  \item \texttt{OptionalityMeter} --- Ch.~25
  \item \texttt{MaxOptionalityAI} --- Ch.~25
  \item \texttt{BranchingSimulator} --- Ch.~26
  \item \texttt{ScenariosDatabase} --- Ch.~26
  \item \texttt{HistoryFunctor} --- Ch.~27
  \item \texttt{VictoryChecker} --- Ch.~27
  \item \texttt{IntelligenceFusion} --- Ch.~28
  \item \texttt{GlobalStateReconstructor} --- Ch.~28
  \item \texttt{LogMerge} --- Ch.~29
  \item \texttt{DiplomaticIntelligibility} --- Ch.~30
  \item \texttt{CulturalMissionary} --- Ch.~30
  \item \texttt{AutopoiesisCheck} --- Ch.~31
  \item \texttt{FixedMul8} --- Ch.~2
  \item \texttt{PredictionErrorTracker} --- Ch.~13
  \item \texttt{BlanketStabilityMeter} --- Ch.~13
  \item \texttt{CollectiveBlanket} --- Ch.~13
\end{itemize}
\end{multicols}

\end{document}